\documentclass[10pt,landscape, a4paper]{article}
\usepackage{anyfontsize}
\renewcommand\normalsize{\fontsize{12}{7}\selectfont}
\usepackage{multicol}
\usepackage{calc}
\usepackage{ifthen}
\usepackage[landscape]{geometry}
\usepackage{fancyhdr}
\usepackage{amsmath,amsthm,amsfonts,amssymb}
\usepackage{color,graphicx,overpic}
\usepackage{hyperref}
\usepackage{listings}
\usepackage{xcolor}
\renewcommand{\ttdefault}{pcr}
\lstset{
    basicstyle=\ttfamily,
    keywordstyle=\color{blue}\bfseries,
    stringstyle=\color{red},
    commentstyle=\color{gray},
    identifierstyle=\color{black},
    columns=fullflexible,
    frame=single,
    breaklines=true,
    showstringspaces=false,
    rulecolor=\color{black},
    framerule=0.4pt
}

\usepackage{titlesec}

\titlespacing*{\section}
{0pt}    % left
{0.3ex}  % space BEFORE
{0.2ex}  % space AFTER


\pdfinfo{
  /Title (example.pdf)
  /Creator (TeX)
  /Producer (pdfTeX 1.40.0)
  /Author (Om Amar)
  /Subject (ICPC TDR)
  /Keywords (pdflatex, latex,pdftex,tex)}

% This sets page margins to .5 inch if using letter paper, and to 1cm
% if using A4 paper. (This probably isn't strictly necessary.)
% If using another size paper, use default 1cm margins.
\ifthenelse{\lengthtest { \paperwidth = 11in}}
    { \geometry{top=0.6in,left=.1in,right=.1in,bottom=.1in} }
    {\ifthenelse{ \lengthtest{ \paperwidth = 297mm}}
        {\geometry{top=1.7cm,left=0.5cm,right=0.5cm,bottom=0.2cm} }
        {\geometry{top=1.7cm,left=0.5cm,right=0.5cm,bottom=0.2cm} }
    }

% Turn off header and footer
\pagestyle{fancy}
\fancyfoot{}

\fancyhead[L]{\Large Team Blank, IIIT Nagpur}
\fancyhead[R]{\Large \thepage}
% Redefine section commands to use less space

% Define BibTeX command
\def\BibTeX{{\rm B\kern-.05em{\sc i\kern-.025em b}\kern-.08em
    T\kern-.1667em\lower.7ex\hbox{E}\kern-.125emX}}

% Don't print section numbers
\setcounter{secnumdepth}{0}


\setlength{\parindent}{0pt}
\setlength{\parskip}{0pt plus 0.5ex}

%My Environments
\newtheorem{example}[section]{Example}
% -----------------------------------------------------------------------



\begin{document}

% ---------- COVER PAGE ----------
\thispagestyle{empty}
\setcounter{page}{0}

\vspace*{\fill}
\begin{center}
{\Huge \textbf{Indian Institute of Information Technology, Nagpur}}\\[2cm]
{\LARGE \textbf{Team Blank}}
\end{center}
\vspace*{\fill}

\clearpage
% ---------- END COVER PAGE ----------

% Reset page layout cleanly
\setcounter{page}{1}
\pagestyle{fancy}

\begin{multicols*}{3}

% multicol parameters
% These lengths are set only within the two main columns
%\setlength{\columnseprule}{0.25pt}
\setlength{\premulticols}{1pt}
\setlength{\postmulticols}{1pt}
\setlength{\multicolsep}{1pt}
\setlength{\columnsep}{2pt}

\section{I/O}
\begin{lstlisting}[language=Python]
import sys,math,cmath,random,os,psutil
from heapq import heappush,heappop
from bisect import bisect_right,bisect_left
from collections import Counter,deque,defaultdict
from itertools import permutations,combinations
from io import BytesIO, IOBase
from decimal import Decimal,getcontext
process = psutil.Process(os.getpid())
BUFSIZE = 8192
class FastIO(IOBase):
    newlines = 0
    def __init__(self, file):
        self._file = file
        self._fd = file.fileno()
        self.buffer = BytesIO()
        self.writable = "x" in file.mode or "r" not in file.mode
        self.write = self.buffer.write if self.writable else None
    def read(self):
        while True:
            b = os.read(self._fd, max(os.fstat(self._fd).st_size, BUFSIZE))
            if not b:
                break
            ptr = self.buffer.tell()
            self.buffer.seek(0, 2), self.buffer.write(b), self.buffer.seek(ptr)
        self.newlines = 0
        return self.buffer.read()
    def readline(self):
        while self.newlines == 0:
            b = os.read(self._fd, max(os.fstat(self._fd).st_size, BUFSIZE))
            self.newlines = b.count(b"\n") + (not b)
            ptr = self.buffer.tell()
            self.buffer.seek(0, 2), self.buffer.write(b), self.buffer.seek(ptr)
        self.newlines -= 1
        return self.buffer.readline()
    def flush(self):
        if self.writable:
            os.write(self._fd, self.buffer.getvalue())
            self.buffer.truncate(0), self.buffer.seek(0)
class IOWrapper(IOBase):
    def __init__(self, file):
        self.buffer = FastIO(file)
        self.flush = self.buffer.flush
        self.writable = self.buffer.writable
        self.write = lambda s: self.buffer.write(s.encode("ascii"))
        self.read = lambda: self.buffer.read().decode("ascii")
        self.readline = lambda: self.buffer.readline().decode("ascii")
sys.stdin, sys.stdout = IOWrapper(sys.stdin), IOWrapper(sys.stdout)
MOD = 10**9 + 7
RANDOM = random.randrange(1,2**62)
def gcd(a,b):
    while b:a,b = b,a%b;return a
def lcm(a,b):return a//gcd(a,b)*b
II = lambda : int(sys.stdin.readline().strip())
LII = lambda : list(map(int, sys.stdin.readline().split()))
SI = lambda : sys.stdin.readline().strip()
LII_1 = lambda : list(map(lambda x:int(x)-1, sys.stdin.readline().split()))
if True:
    base = os.path.dirname(os.path.abspath(__file__))
    sys.stdin  = open(os.path.join(base, "input.txt"), "r")
    sys.stdout = open(os.path.join(base, "output.txt"), "w")
    sys.stderr = open(os.path.join(base, "error.txt"), "w")
def solve():
    return
solve()
\end{lstlisting}

\section{Number Theory}
\begin{lstlisting}[language=Python]
def miller_is_prime(n):
    if n < 5 or n & 1 == 0 or n % 3 == 0:return 2 <= n <= 3
    s = ((n - 1) & (1 - n)).bit_length() - 1
    d = n >> s
    for a in [2, 325, 9375, 28178, 450775, 9780504, 1795265022]:
        p = pow(a, d, n)
        if p == 1 or p == n - 1 or a % n == 0:continue
        for _ in range(s):
            p = (p * p) % n
            if p == n - 1:break
        else:return False
    return True
def is_prime(n):
    if n <= 1:return False
    if n <= 3:return True
    if n % 2 == 0 or n % 3 == 0:return False
    i = 5
    while i * i <= n:
        if n % i == 0 or n % (i + 2) == 0:return False
        i += 6
    return True
\end{lstlisting}
\begin{lstlisting}[language=Python]
def sieve(n):
    primes = [];isp = [1] * (n+1)
    isp[0] = isp[1] = 0
    for i in range(2,n+1):
        if isp[i]:
            primes.append(i)
            for j in range(i*i,n+1,i):isp[j] = 0
    return primes
def sieve_unique(N):
    mini = [i for i in range(N)]
    for i in range(2,N):
        if mini[i]==i:
            for j in range(2*i,N,i):
                mini[j] = i
    return mini
MAX_N = 10**6+1
Lmini = sieve_unique(MAX_N)
def prime_factors(k,typ=0):
    if typ==0:ans = Counter()
    elif typ==1:ans = set()
    else:ans = []
    while k!=1:
        if typ==0:ans[Lmini[k]] += 1
        elif typ==1:ans.add(Lmini[k])
        else:ans.append(Lmini[k])
        k //= Lmini[k]
    return ans
def all_factors(x):
    L = list(prime_factors(x).items())
    st = [1]
    for i in range(len(L)):
        for j in range(len(st)-1,-1,-1):
            k = L[i][0]
            for l in range(L[i][1]):
                st.append(st[j]*k)
                k *= L[i][0]
    return st
\end{lstlisting}
\begin{lstlisting}[language=Python]

def pollard_rho(n):
    # works upto 1e18
    if (n == 1):
        return n
    if (n % 2 == 0):
        return 2
    for _ in range(50):
        x = random.randint(2,n-1)
        c = random.randint(1,n-1)
        y = x
        d = 1
        while (d == 1):
            x = (pow(x, 2, n) + c + n)%n
            y = (pow(y, 2, n) + c + n)%n
            y = (pow(y, 2, n) + c + n)%n
            d = gcd(abs(x - y), n)
            if (d == n):
                break
        if d!=1 and d!=n:
            return n

    return None

def miller_pollard_rho(n):
    # returns a random factor of n, works upto 1e18
    if n&1==0:return 2
    if n%3==0:return 3
    s = ((n-1)&(1-n)).bit_length() - 1
    d = n>>s
    for a in [2, 325, 9375, 28178, 450775, 9780504, 1795265022]:
        p = pow(a, d, n)
        if p == 1 or p == n - 1 or a % n == 0:continue
        for _ in range(s):
            prev = p
            p = (p * p) % n
            if p == 1:return gcd(prev - 1, n)
            if p == n - 1:break
        else:
            for i in range(2, n):
                x, y = i, (i * i + 1) % n
                f = gcd(abs(x - y), n)
                while f == 1:
                    x, y = (x * x + 1) % n, (y * y + 1) % n
                    y = (y * y + 1) % n
                    f = gcd(abs(x - y), n)
                if f != n:return f
    return n

@memodict
def prime_factors_large(n):
    # returns prime factor in n^(1/4) but is probablistic
    if n <= 1:return Counter()
    f = miller_pollard_rho(n) #pollard_rho is another option
    return Counter([n]) if f == n else prime_factors_large(f) + prime_factors_large(n // f)

def discrete_log(a, b, mod):
    # returns smalest x such that pow(a,x,mod) = b
    # given gcd(a,mod)==1
    n = int(mod**0.5) + 1
    tiny_step, e = {}, 1
    for j in range(1, n + 1):
        e = e * a % mod
        if e == b:
            return j
        tiny_step[b * e % mod] = j
    factor = e
    for i in range(2, n + 2):
        e = e * factor % mod
        if e in tiny_step:
            j = tiny_step[e]
            return n * i - j if pow(a, n * i - j, mod) == b else None
    return None

def extended_gcd(a, b):
    # returns gcd(a, b), s, r s.t. a * s + b * r == gcd(a, b)
    s, old_s = 0, 1
    r, old_r = b, a
    while r:
        q = old_r // r
        old_r, r = r, old_r - q * r
        old_s, s = s, old_s - q * s
    return old_r, old_s, (old_r - old_s * a) // b if b else 0

def composite_crt(b, m):
    # returns x s.t. x = b[i] (mod m[i]) for all i
    x, m_prod = 0, 1
    for bi, mi in zip(b, m):
        g, s, _ = extended_gcd(m_prod, mi)
        if ((bi - x) % mi) % g:
            return None
        x += m_prod * (s * ((bi - x) % mi) // g)
        m_prod = (m_prod * mi) // gcd(m_prod, mi)
    return x % m_prod

def phi(n):
    ph = [i if i & 1 else i // 2 for i in range(n + 1)]
    for i in range(3,n+1,2):
        if ph[i]==i:
            for j in range(i,n+1,i):
                ph[j] = (ph[j]//i)*(i-1)
    return ph
\end{lstlisting}


\section{Matrix}
\begin{lstlisting}[language=Python]
class Matrix:
    MOD = 10**9+7
    @staticmethod
    def transpose(mat):
        return [list(c) for c in zip(*mat)]
    def __init__(self,M):
        self.M=M
    def minmul(self,o):
        L,B=self.M,o.M
        r=len(L); c=len(B[0])
        return Matrix([[min(L[i][k]+B[k][j] for k in range(len(B)))%self.MOD for j in range(c)] for i in range(r)])
    def __matmul__(self,o):
        L,B=self.M,o.M
        r=len(L); c=len(B[0])
        return Matrix([[sum(L[i][k]*B[k][j] for k in range(len(B)))%self.MOD for j in range(c)] for i in range(r)])
    def __pow__(self,p):
        M=self; n=len(M.M)
        R=Matrix([[i==j for j in range(n)] for i in range(n)])
        while p:
            R = R@M if p&1 else R
            M = M@M
            p//=2
        return R
    @staticmethod
    def gauss(A,mod=10**9+7):
        m,n=len(A),len(A[0])-1; r=0; L=[-1]*n
        for c in range(n):
            for i in range(r,m):
                if A[i][c]: A[r],A[i]=A[i],A[r]
                break
            else:
                continue
            k=pow(A[r][c],-1,mod)
            for j in range(c,n+1):
                A[r][j]=A[r][j]*k%mod
            for i in range(m):
                if i!=r and A[i][c]:
                    f=A[i][c]
                    for j in range(c,n+1):
                        A[i][j]=(A[i][j]-f*A[r][j])%mod
            L[c]=r
            r+=1
        if any(A[i][n] for i in range(r,m)):
            return None
        return [A[L[i]][n] if L[i]!=-1 else 0 for i in range(n)]
    def __str__(self):
        return '\n'.join(' '.join(map(str,row)) for row in self.M)
    def __repr__(self):
        return f"Matrix({self.M})"
    def __iter__(self):
        return iter(self.M)
\end{lstlisting}

\section{Combinatorics}
\begin{lstlisting}[language=Python]

class Factorial:
    def __init__(self, N, mod):
        self.mod = mod
        self.f = [1 for _ in range(N)]
        self.g = [1 for _ in range(N)]
        for i in range(1, N):self.f[i] = self.f[i - 1] * i % self.mod
        self.g[-1] = pow(self.f[-1], mod - 2, mod)
        for i in range(N - 2, -1, -1):self.g[i] = self.g[i + 1] * (i + 1) % self.mod
    def derangements(self,n):
        # Number of permutations of n objects where no object appears in its original position
        return int(self.fac(n)/math.e+0.5)
    def der(self, n):
        if n==0: return 1
        if n==1: return 0
        a, b = 1, 0  # der(0), der(1)
        for i in range(2,n+1):
            a,b = b,((i-1)*(a+b)) % self.mod
        return b
    def stirling_2(self,n,k):
        # Number of ways to partition n elements into k non-empty subsets
        return sum(((-1)**(k - j)) * math.comb(k, j) * (j**n) for j in range(k + 1)) // math.factorial(k)
    def stirling_2_mod(self,n,k):
        return (sum(((-1)**(k - j))*self.combi(k, j)*pow(j,n,self.mod) for j in range(k + 1)) * self.fac_inv(k))%self.mod
    def partition(self,n,k):
        # Ways to partition n into k or fewer parts of size 1 or greater
        # add dp[(n,k)] and memoize it if using it.
        if n<0 or (k<1 and n>0):return 0
        if n==0:return 1
        return 1 if n==1 else self.partition(n,k-1)+self.partition(n-k,k)
    def lucas_nCk(self,n,k):
        # Lucas's theorem for finding ((nCk)%p) for large n,k
        # in log p(n)
        ans = 1
        p = self.mod
        while n or k:
            x = n%p
            y = k%p
            ans *= self.combi(x,y)
            ans %= p
            n //= p
            k //= p
        return ans
    def catalan(self, n):
        return (self.combi(2 * n, n) - self.combi(2 * n, n - 1)) % self.mod
    def fac(self, n):
        return self.f[n]
    def fac_inv(self, n):
        return self.g[n]
    def combi(self, n, m):
        if m == 0: return 1
        if n < m or m < 0 or n < 0: return 0
        return self.f[n] * self.g[m] % self.mod * self.g[n - m] % self.mod
    def permu(self, n, m):
        if n < m or m < 0 or n < 0: return 0
        return self.f[n] * self.g[n - m] % self.mod
    def inv(self, n):
        return self.f[n-1] * self.g[n] % self.mod
    def beggar(self,n,sumi):
        # number of ways n elements >=0 have sum sumi
        return self.combi(n+sumi-1,sumi)
\end{lstlisting}

\section{Geometry}
\begin{lstlisting}[language=Python]
## here onwards, it is for intersection only, even the point class is different
class Point:
    def __init__(self, x, y):
        self.x = x
        self.y = y
def onSegment(p, q, r):
    if ( (q.x <= max(p.x, r.x)) and (q.x >= min(p.x, r.x)) and
           (q.y <= max(p.y, r.y)) and (q.y >= min(p.y, r.y))):
        return True
    return False
def orientation(p, q, r):
    # to find the orientation of an ordered triplet (p,q,r)
    # function returns the following values:
    # 0 : Collinear points
    # 1 : Clockwise points
    # 2 : Counterclockwise
    val = ((q.y - p.y) * (r.x - q.x)) - ((q.x - p.x) * (r.y - q.y))
    if (val > 0):
        return 1
    elif (val < 0):
        return 2
    else:
        return 0
def doIntersect(p1,q1,p2,q2):
    o1 = orientation(p1, q1, p2)
    o2 = orientation(p1, q1, q2)
    o3 = orientation(p2, q2, p1)
    o4 = orientation(p2, q2, q1)
    if ((o1 != o2) and (o3 != o4)):
        return True
    if ((o1 == 0) and onSegment(p1, p2, q1)):
        return True
    if ((o2 == 0) and onSegment(p1, q2, q1)):
        return True
    if ((o3 == 0) and onSegment(p2, p1, q2)):
        return True
    if ((o4 == 0) and onSegment(p2, q1, q2)):
        return True
    return False
##
## lines
# 2d line: ax + by + c = 0  is  (a, b, c)
#		  ax + by + c = 0	 ((a, b, c),
# 3d line: dx + ez + f = 0  is  (d, e, f),
#		  gy + hz + i = 0	  (g, h, i))
def get_2dline(p1, p2):
    if p1 == p2:
        return (0, 0, 0)
    _p1, _p2 = min(p1, p2), max(p1, p2)
    a, b, c = _p2[1] - _p1[1], _p1[0] - _p2[0], _p1[1] * _p2[0] - _p1[0] * _p2[1]
    g = gcd(gcd(a, b), c)
    return (a // g, b // g, c // g)
dist = lambda p1, p2: sum((a - b) * (a - b) for a, b in zip(p1, p2))**0.5
get_line = lambda p1, p2: map(get_2dline, combinations(p1, 2), combinations(p2, 2))
is_parallel = lambda l1, l2: l1[0] * l2[1] == l2[0] * l1[1]
is_same = lambda l1, l2: is_parallel(l1, l2) and (l1[1] * l2[2] == l2[1] * l1[2])
collinear = lambda p1, p2, p3: is_same(get_2dline(p1, p2), get_2dline(p2, p3))
intersect = (lambda l1, l2: None if is_parallel(l1, l2) else (
    (l2[1] * l1[2] - l1[1] * l2[2]) / (l2[0] * l1[1] - l1[0] * l2[1]),
    (l1[0] * l2[2] - l1[2] * l2[0]) / (l2[0] * l1[1] - l1[0] * l2[1]),
))
rotate = lambda p, theta, origin=(0, 0): (
    origin[0] + (p[0] - origin[0]) * math.cos(theta) - (p[1] - origin[1]) * math.sin(theta),
    origin[1] + (p[0] - origin[0]) * math.sin(theta) + (p[1] - origin[1]) * math.cos(theta),
)
##
## polygons
dist = lambda p1, p2: sum((a - b) * (a - b) for a, b in zip(p1, p2))**0.5
perimeter = lambda *p: sum(dist(i, j) for i, j in zip(p, p[1:] + p[:1]))
# shoelace theorem
area = lambda *p: abs(sum(i[0] * j[1] - j[0] * i[1] for i, j in zip(p, p[1:] + p[:1]))) / 2
centroid = lambda p: (
    sum((i[0] + j[0]) * (i[0] * j[1] - j[0] * i[1])
        for i, j in zip(p, p[1:] + p[:1]))
    / (3 * sum(i[0] * j[1] - j[0] * i[1] for i, j in zip(p, p[1:] + p[:1]))),
    sum((i[1] + j[1]) * (i[0] * j[1] - j[0] * i[1])
        for i, j in zip(p, p[1:] + p[:1]))
    / (3 * sum(i[0] * j[1] - j[0] * i[1] for i, j in zip(p, p[1:] + p[:1])))
)
## above 2 functions can be done on a range of points in O(1)
is_in_circle = lambda p, c, r: sum(i * i - j * j for i, j in zip(p, c)) < r * r
incircle_radius = lambda a, b, c: area(a, b, c) / (perimeter(a, b, c) / 2)
circumcircle_radius = lambda a, b, c: (dist(a, b) * dist(b, c) * dist(c, a)) / (4 * area(a, b, c))
##
def convex_hull(P):
    if len(P)<=1:return P
    cross = lambda a,b,c: (a[0] - b[0]) * (c[1] - b[1]) - (a[1] - b[1]) * (c[0] - b[0])
    P = sorted(P)
    def f(P):
        lower = []
        for p in P:
            while len(lower) >= 2 and cross(lower[-2], lower[-1], p) < 0:
                lower.pop()
            lower.append(p)
        return lower
    return f(P)[:-1]+f(P[::-1])[:-1]
\end{lstlisting}

\section{Linear Algebra}
\begin{lstlisting}[language=Python]

def solve_quadratic(a, b, c):
    if a == 0:return None
    D = b*b - 4*a*c
    if D < 0:return None
    s = math.isqrt(D)
    if s*s != D:return None
    den = 2*a
    if den == 0:return None
    x1, r1 = divmod(-b + s, den)
    x2, r2 = divmod(-b - s, den)
    if r1!=0 or r2!=0:return None
    return (x1, x2)

def solve_cubic(a, b, c, d):
    if a == 0:return None
    roots = []
    for x in range(-math.isqrt(abs(d)) - 1, math.isqrt(abs(d)) + 2):
        if a*x*x*x + b*x*x + c*x + d == 0:
            roots.append(x)
    return tuple(roots) if roots else None

def compute_rank(A):
    EPS = 1e-9
    A = [row[:] for row in A]
    n = len(A)
    m = len(A[0])
    rank = 0
    row_selected = [False] * n
    for i in range(m):
        j = 0
        while j < n:
            if not row_selected[j] and abs(A[j][i]) > EPS:
                break
            j += 1
        if j != n:
            rank += 1
            row_selected[j] = True
            for p in range(i + 1, m):
                A[j][p] /= A[j][i]
            for k in range(n):
                if k != j and abs(A[k][i]) > EPS:
                    for p in range(i + 1, m):
                        A[k][p] -= A[j][p] * A[k][i]
    return rank

class XorBasis:
    def __init__(self,A=[]):
        self.reduced_base = {}
        self.how = {}
        self.base = {}
        for i in A:
            self.add(i)
    def max(self):
        base = self.base; reduced_base = self.reduced_base
        how = self.how
        x = 0;tmp = 0;ans=[]
        for j in sorted(reduced_base, reverse=True):
            if not x & (1 << j):
                x ^= reduced_base[j]
                tmp ^= how[j]
                ans.append(base[j])
        # xor of these elements return the max element
        return ans
    def make(self,k):
        # returns True if can make x with these elements
        reduced_base = self.reduced_base
        for j in range(k.bit_length()-1,-1,-1):
            if ((1<<j)&k):
                if j in reduced_base:
                    k ^= reduced_base[j]
                else:
                    return False
        return k==0
    def add(self,a):
        # adds a to the basis
        reduced_base = self.reduced_base
        how = self.how;base = self.base
        tmp = 0
        while a:
            b = a.bit_length() - 1
            if b in reduced_base:
                a ^= reduced_base[b]
                tmp ^= how[b]
            else:
                reduced_base[b] = a
                how[b] = tmp | (1 << len(base))
                base[b] = a
                break

def simplex_bland(c, A, b, frac=False):
    # maximize ci.xi
    # given ai.xi <= bi and xi >= 0
    if not frac:
        n = len(c)
        m = len(b)
        T = [A[i] + [0]*i + [1] + [0]*(m-i-1) + [b[i]] for i in range(m)]
        T.append([-ci for ci in c] + [0]*m + [0])
        N = list(range(n)) # non-basic
        B = list(range(n, n+m)) # basic
        def pivot(r, s):
            # Make coeffecient of pivot row 1 and Eliminate column s in all other rows
            T[r] = [v / T[r][s] for v in T[r]]
            for i in range(len(T)):
                if i != r:
                    fac = T[i][s]
                    T[i] = [T[i][j] - fac * T[r][j] for j in range(len(T[0]))]
        while True:
            s = None
            for j in range(len(T[0]) - 1):
                if T[-1][j] < 0:
                    s = j
                    break
            if s is None:
                break
            # leaving row
            r = None
            for i in range(m):
                if T[i][s] > 1e-12:
                    ratio = T[i][-1] / T[i][s]
                    if r is None or ratio < T[r][-1] / T[r][s]:
                        r = i
            if r is None:
                return None
            pivot(r, s)
            B[r], N[s] = N[s], B[r]
        x = [0] * (n + m)
        for i in range(m):
            x[B[i]] = T[i][-1]
        return x[:n], T[-1][-1]
    else:
        from fractions import Fraction
        n = len(c)
        m = len(b)
        c = [Fraction(ci) for ci in c]
        A = [[Fraction(aij) for aij in row] for row in A]
        b = [Fraction(bi) for bi in b]
        T = [A[i] + [Fraction(1 if i == j else 0) for j in range(m)] + [b[i]]
             for i in range(m)]
        T.append([-ci for ci in c] + [Fraction(0)] * (m + 1))
        N = list(range(n))
        B = list(range(n, n+m))
        def pivot(r, s):
            piv = T[r][s]
            T[r] = [v / piv for v in T[r]]
            for i in range(m + 1):
                if i != r:
                    fac = T[i][s]
                    T[i] = [T[i][j] - fac * T[r][j] for j in range(len(T[0]))]
        while True:
            s = None
            for j in range(len(T[0]) - 1):
                if T[-1][j] < 0:
                    s = j
                    break
            if s is None:
                break
            r = None
            for i in range(m):
                if T[i][s] > 0:
                    ratio = T[i][-1] / T[i][s]
                    if r is None or ratio < T[r][-1] / T[r][s]:
                        r = i
            if r is None:
                return None
            pivot(r, s)
            B[r], N[s] = N[s], B[r]
        x = [Fraction(0)] * (n + m)
        for i in range(m):
            x[B[i]] = T[i][-1]
        return x[:n], T[-1][-1]

def s_lr(l,r,pow=1):
    if l>r:return 0
    if pow==1:return (r*(r+1)//2) - ((l-1)*l//2)
    elif pow==2:return (r*(r+1)*(2*r+1)//6) - ((l-1)*l*(2*l-1)//6)
    elif pow==3:return (r*(r+1)//2)**2 - ((l-1)*l//2)**2
    else:return None

class DiophantineEquations:
    """
        used for solving equations of the form a*x + b*y = c,
        solnll takes the lower limit as well
    """
    def __init__(self):pass
    def euclidean_gcd(self,a, b):
        """
            euclidean gcd , returns x and y such that
            a*x + b*y = gcd(a,b)
        """
        if b == 0:
            return a,1,0
        g,x1,y1 = self.euclidean_gcd(b, a % b)
        x = y1
        y = x1-(a//b)*y1
        return g,x,y
    def soln(self,a,b,c,t=10**18,t1=10**18):
        """
            return m and n such that a*m + b*n = c and 0<=m<=t and 0<=n<=t1
            don't input t,t1 for any possible value
        """
        g = gcd(a, b)
        if c%g!=0:
            return -1,-1
        _, x, y = self.euclidean_gcd(a, b)
        k3 = x*(c//g)
        n1 = y*(c//g)
        k1 = max(math.ceil(-k3*g/b),math.ceil((n1-t1)/(a/g)))
        kmaxi = min(math.floor((t-k3)/(b/g)), math.floor(n1/(a/g)))
        if k1<=kmaxi:
            k = k1
            m = k3+k*(b//g)
            n = n1-k*(a//g)
            return m, n
        else:
            return -1,-1
    def sollnll(self,a,b,c,t,t1,m_lower,n_lower):
        """
            return m and n such that a*m + b*n = c and
            m_lower <= m <= t and n_lower <= n <= t1
        """
        g = gcd(a, b)
        if c % g != 0:
            return -1, -1
        _, x, y = self.euclidean_gcd(a, b)
        k3 = x * (c // g)
        n1 = y * (c // g)
        k1 = max(math.ceil((m_lower - k3 * g) / b), math.ceil((n1 - t1) / (a / g)))
        kmaxi = min(math.floor((t - k3) / (b / g)), math.floor((n1 - n_lower) / (a / g)))
        if k1 <= kmaxi:
            k = k1
            m = k3 + k * (b // g)
            n = n1 - k * (a // g)
            if m_lower <= m <= t and n_lower <= n <= t1:
                return m, n
            else:
                return -1,-1
        else:
            return -1, -1
\end{lstlisting}

\section{Data Structures}
Disjoint Set Union
\begin{lstlisting}[language=Python]
class DisjointSetUnion:
    def __init__(self, n):
        self.parent = list(range(n))
        # self.s = set(self.parent)
        self.size = [1] * n
    def find(self, a):
        acopy = a
        while a != self.parent[a]:
            a = self.parent[a]
        while acopy != a:
            self.parent[acopy], acopy = a, self.parent[acopy]
        return a
    def union(self, a, b):
        a, b = self.find(a), self.find(b)
        if a != b:
            if self.size[a] < self.size[b]:
                a, b = b, a
            # self.s.remove(b)
            self.parent[b] = a
            self.size[a] += self.size[b]
    def set_size(self, a):
        return self.size[self.find(a)]
    def __len__(self):
        return len(self.s)
    def notfind(self, a):
        k = self.find(a)
        for j in self.s:
            if j!=k:
                return j
        return -1
\end{lstlisting}

Persistent DSU
\begin{lstlisting}[language=Python]
class PersistentDSU:
    def __init__(self,n):
        self.parent = list(range(n))
        self.size = [1]*n
        self.time = [float('inf')]*n
    def find(self,node,version):
        # returns root at given version
        while not (self.parent[node]==node or self.time[node]>version):
            node = self.parent[node]
        return node
    def union(self,a,b,time):
        # merges a and b
        a = self.find(a,time)
        b = self.find(b,time)
        if a==b:
            return False
        if self.size[a]>self.size[b]:
            a,b = b,a
        self.parent[a] = b
        self.time[a] = time
        self.size[b] += self.size[a]
        return True
    def isconnected(self,a,b,time):
        return self.find(a,time)==self.find(b,time)

        
class RollbackDSU:
    def __init__(self,n):
        self.p=list(range(n)); self.sz=[1]*n; self.st=[]
    def find(self,x):
        while x!=self.p[x]: x=self.p[x]
        return x
    def union(self,a,b):
        a,b=self.find(a),self.find(b)
        if a==b: self.st.append((-1,-1,-1)); return
        if self.sz[a]>self.sz[b]: a,b=b,a
        self.st.append((a,b,self.sz[b]))
        self.p[a]=b; self.sz[b]+=self.sz[a]
    def rollback(self):
        a,b,s=self.st.pop()
        if a==-1: return
        self.p[a]=a; self.sz[b]=s
\end{lstlisting}

Fenwick Tree
\begin{lstlisting}[language=Python]
class BIT:
    #Faster than segment tree so use if possbile
    def __init__(self, x):
        self.bit = x
        for i in range(len(x)):
            j = i | (i + 1)
            if j < len(x):
                x[j] += x[i]
    def update(self, idx, x):
        """updates bit[idx] += x"""
        while idx < len(self.bit):
            self.bit[idx] += x
            idx |= idx + 1
    def query(self, end):
        """calc sum(bit[:end])"""
        #gives sum of element before end
        x = 0
        while end:
            x += self.bit[end - 1]
            end &= end - 1
        return x
    def findkth(self, k):
        """Find largest idx such that sum(bit[:idx]) <= k"""
        idx = -1
        for d in reversed(range(len(self.bit).bit_length())):
            right_idx = idx + (1 << d)
            if right_idx < len(self.bit) and k >= self.bit[right_idx]:
                idx = right_idx
                k -= self.bit[idx]
        return idx + 1
\end{lstlisting}

2-D Fenwick Tree
\begin{lstlisting}[language=Python]
class BIT2D:
    def __init__(self, arr):
        self.n = len(arr)
        self.m = len(arr[0]) if self.n > 0 else 0
        # self.bit = [row[:] for row in arr]
        self.bit = arr # assuming that arr is not used after this
        for i in range(self.n):
            for j in range(self.m):
                ni = i | (i + 1)
                if ni < self.n:
                    self.bit[ni][j] += self.bit[i][j]
        for i in range(self.n):
            for j in range(self.m):
                nj = j | (j + 1)
                if nj < self.m:
                    self.bit[i][nj] += self.bit[i][j]
    def add(self, x, y, delta):
        # 0-based in log n * log m
        i = x
        while i < self.n:
            j = y
            while j < self.m:
                self.bit[i][j] += delta
                j |= j + 1
            i |= i + 1
    def sum(self, x, y):
        # sum from 0,0 to x,y inclusive in log n * log m
        if not (0<=x<self.n) or not (0<=y<self.m):
            return 0
        res = 0
        i = x
        while i >= 0:
            j = y
            while j >= 0:
                res += self.bit[i][j]
                j = (j & (j + 1)) - 1
            i = (i & (i + 1)) - 1
        return res
    def query(self, x1, y1, x2, y2):
        # sum of L[x1:x2+1][y1:y2+1]
        return (self.sum(x2,y2)-self.sum(x1-1,y2)-self.sum(x2,y1-1)+(self.sum(x1-1,y1-1)))
\end{lstlisting}

Bucket Sorted List
\begin{lstlisting}[language=Python]
from typing import Generic, Iterable, Iterator, List, Tuple, TypeVar, Optional
T = TypeVar('T')
class SortedList(Generic[T]):
    BUCKET_RATIO = 16
    SPLIT_RATIO = 24
    def __init__(self, a: Iterable[T] = []) -> None:
        a = list(a)
        n = self.size = len(a)
        if any(a[i] > a[i + 1] for i in range(n - 1)):
            a.sort()
        num_bucket = int(math.ceil(math.sqrt(n / self.BUCKET_RATIO)))
        self.a = [a[n * i // num_bucket : n * (i + 1) // num_bucket] for i in range(num_bucket)]
    def __iter__(self) -> Iterator[T]:
        for i in self.a:
            for j in i: yield j
    def __reversed__(self) -> Iterator[T]:
        for i in reversed(self.a):
            for j in reversed(i): yield j
    def __eq__(self, other) -> bool:return list(self) == list(other)
    def __len__(self) -> int:return self.size
    def __repr__(self) -> str:return "SortedMultiset" + str(self.a)
    def __str__(self) -> str:
        s = str(list(self))
        return "{" + s[1 : len(s) - 1] + "}"
    def _position(self, x: T) -> Tuple[List[T], int, int]:
        for i, a in enumerate(self.a):
            if x <= a[-1]: break
        return (a, i, bisect_left(a, x))
    def __contains__(self, x: T) -> bool:
        if self.size == 0: return False
        a, _, i = self._position(x)
        return i != len(a) and a[i] == x
    def count(self, x: T) -> int:return self.index_right(x) - self.index(x)
    def add(self,x):return self.insert(x)
    def insert(self, x: T) -> None:
        if self.size == 0:
            self.a = [[x]]
            self.size = 1
            return
        a, b, i = self._position(x)
        a.insert(i, x)
        self.size += 1
        if len(a) > len(self.a) * self.SPLIT_RATIO:
            mid = len(a) >> 1
            self.a[b:b+1] = [a[:mid], a[mid:]]
    def _pop(self, a: List[T], b: int, i: int) -> T:
        ans = a.pop(i)
        self.size -= 1
        if not a: del self.a[b]
        return ans
    def remove(self, x: T) -> bool:
        if self.size == 0: return False
        a, b, i = self._position(x)
        if i == len(a) or a[i] != x: return False
        self._pop(a, b, i)
        return True
    def lt(self, x: T) -> Optional[T]:
        for a in reversed(self.a):
            if a[0] < x:
                return a[bisect_left(a, x) - 1]
    def le(self, x: T) -> Optional[T]:
        for a in reversed(self.a):
            if a[0] <= x:
                return a[bisect_right(a, x) - 1]
    def gt(self, x: T) -> Optional[T]:
        for a in self.a:
            if a[-1] > x:
                return a[bisect_right(a, x)]
    def ge(self, x: T) -> Optional[T]:
        for a in self.a:
            if a[-1] >= x:
                return a[bisect_left(a, x)]
    def __getitem__(self, i: int) -> T:
        if i < 0:
            for a in reversed(self.a):
                i += len(a)
                if i >= 0: return a[i]
        else:
            for a in self.a:
                if i < len(a): return a[i]
                i -= len(a)
        raise IndexError
    def pop(self, i: int = -1) -> T:
        if i < 0:
            for b, a in enumerate(reversed(self.a)):
                i += len(a)
                if i >= 0: return self._pop(a, ~b, i)
        else:
            for b, a in enumerate(self.a):
                if i < len(a): return self._pop(a, b, i)
                i -= len(a)
        raise IndexError
    def bisect_left(self,x):return self.index(x)
    def index(self, x: T) -> int:
        ans = 0
        for a in self.a:
            if a[-1] >= x:
                return ans + bisect_left(a, x)
            ans += len(a)
        return ans
    def bisect_right(self,x):return self.index_right(x)
    def index_right(self, x: T) -> int:
        ans = 0
        for a in self.a:
            if a[-1] > x:
                return ans + bisect_right(a, x)
            ans += len(a)
        return ans
    def find_closest(self, k: T) -> Optional[T]:
        if self.size == 0:return None
        ltk = self.le(k);gtk = self.ge(k)
        if ltk is None:return gtk
        if gtk is None:return ltk
        return ltk if abs(k-ltk)<=abs(k-gtk) else gtk
\end{lstlisting}

Trie
\begin{lstlisting}[language=Python]
class Trie:
    def __init__(self):
        self.root = {}
    def add(self, word):
        current_dict = self.root
        for letter in word:
            current_dict = current_dict.setdefault(letter, {})
        current_dict[0] = True
\end{lstlisting}

Convex Hull
\begin{lstlisting}[language=Python]
class LiChao:
    __slots__ = ("lo","hi","mid","m","b","left","right")
    # lo = min value of x that you will query
    def __init__(self,lo,hi):
        self.lo = lo
        self.hi = hi
        self.mid = (lo+hi)>>1
        self.m = None
        self.b = None
        self.left = None
        self.right = None
    def add_line(self,m,b):
        node = self
        while True:
            lo,hi,mid = node.lo,node.hi,node.mid
            if node.m is None:
                node.m,node.b = m,b
                return
            cur_m,cur_b = node.m,node.b
            if m*mid+b>cur_m*mid+cur_b:
                node.m,m = m,cur_m
                node.b,b = b,cur_b
                cur_m,cur_b = node.m,node.b
            if lo==hi:
                return
            if m*lo+b>cur_m*lo+cur_b:
                if node.left is None:
                    node.left = LiChao(lo,mid)
                node = node.left
            elif m*hi+b>cur_m*hi +cur_b:
                if node.right is None:
                    node.right = LiChao(mid+1,hi)
                node = node.right
            else:
                return
    def query(self,x):
        node = self
        res = -10**18
        while node:
            if node.m is not None:
                v = node.m*x+node.b
                if v>res:
                    res = v
            if node.lo==node.hi:
                break
            if x<=node.mid:
                node = node.left
            else:
                node = node.right
        return res
\end{lstlisting}

Heavy Light Decomposition
\begin{lstlisting}[language=Python]
class HLD:
    def __init__(self, adj, values, root=0,func=max,unit=float('-inf')):
        self.adj = adj
        self.values = values
        self.parent = [-1] * len(adj)
        self.depth = [0] * len(adj)
        self.size = [0] * len(adj)
        self.heavy = [-1] * len(adj)
        self.head = [0] * len(adj)
        self.pos = [0] * len(adj)
        self.flat = [0] * len(adj)
        self.unit = unit
        self.func = func
        self._dfs(root)
        self._decompose(root)
        self.seg = SegmentTree([self.values[self.flat[i]] for i in range(len(self.adj))],func,unit)
    def _dfs(self,start=0):
        visited = [False] * len(self.adj)
        stack = [start]
        while stack:
            start = stack[-1]
            if not visited[start]:
                visited[start] = True
                for child in self.adj[start]:
                    if not visited[child]:
                        self.parent[child] = start
                        self.depth[child] = self.depth[start]+1
                        stack.append(child)
            else:
                self.size[stack.pop()] = 1
                k = 0
                for child in self.adj[start]:
                    if self.parent[start]!=child:
                        self.size[start] += self.size[child]
                        if self.size[child]>k:
                            k = self.size[child]
                            self.heavy[start] = child
        return visited
    def _decompose(self, root):
        stack = [(root,root)]
        time = 0
        while stack:
            u, h = stack.pop()
            self.head[u] = h
            self.flat[time] = u
            self.pos[u] = time
            time += 1
            for v in reversed(self.adj[u]):
                if v!=self.parent[u] and v!=self.heavy[u]:
                    stack.append((v, v))
            if self.heavy[u] != -1:
                stack.append((self.heavy[u], h))
    def query(self, u, v):
        res = self.unit
        while self.head[u] != self.head[v]:
            if self.depth[self.head[u]] < self.depth[self.head[v]]:
                u, v = v, u
            res = self.func(res, self.seg.query(self.pos[self.head[u]], self.pos[u] + 1))
            u = self.parent[self.head[u]]
        if self.depth[u] > self.depth[v]:
            u, v = v, u
        return self.func(res, self.seg.query(self.pos[u], self.pos[v] + 1))
    def update(self, u, value):
        self.seg.update(self.pos[u], value)
    def update_path(self, u, v, value):
        while self.head[u] != self.head[v]:
            if self.depth[self.head[u]] < self.depth[self.head[v]]:
                u, v = v, u
            self.seg.range_update(self.pos[self.head[u]], self.pos[u], value)
            u = self.parent[self.head[u]]
        if self.depth[u] > self.depth[v]:
            u, v = v, u
        self.seg.range_update(self.pos[u], self.pos[v], value)
    def add_to_subtree(self, u, value):
        self.seg.range_update(self.pos[u], self.pos[u] + self.size[u] - 1, value)
\end{lstlisting}

Mono Deque
\begin{lstlisting}[language=Python]
from typing import Callable, TypeVar, Generic, List
T = TypeVar('T')
class MonoStack(Generic[T]):
    def __init__(self, op: Callable[[T, T], T], e: Callable[[], T]):
        self.s: List[T] = []
        self.sMono: List[T] = []
        self.op = op
        self.e = e
    def push(self, x: T):
        self.s.append(x)
        if not self.sMono:
            self.sMono.append(x)
        else:
            self.sMono.append(self.op(self.sMono[-1], x))
    def pop(self) -> T:
        if not self.s:
            return self.e()
        self.sMono.pop()
        return self.s.pop()
    def get(self) -> T:
        return self.sMono[-1] if self.sMono else self.e()
    def empty(self) -> bool:
        return not self.s
class MonoDeque(Generic[T]):
    def __init__(self, op: Callable[[T, T], T] = lambda a,b: max(a,b), e: Callable[[], T] = lambda : -float('inf')):
        # e is the unit value in form of lambda : 0
        # The Function op must be associative, f(a,b)==f(b,a)
        self.op = op
        self.e = e
        self.front = MonoStack(op, e)
        self.back = MonoStack(op, e)
    def push_back(self, x: T):
        self.back.push(x)
    def pop_front(self):
        if self.front.empty():
            while not self.back.empty():
                self.front.push(self.back.pop())
        self.front.pop()
    def get(self) -> T:
        # returns the function op of the elements present in the deque
        return self.op(self.front.get(), self.back.get())
\end{lstlisting}

Merge Sort Tree
\begin{lstlisting}[language=Python]
class MergeSortTree:
    def merge(L,L1):
        L2 = []
        L = list(L);L1 = list(L1)
        i = j = 0
        while j<len(L) and i<len(L1):
            if L[j]<L1[i]:
                L2.append(L[j])
                j += 1
            else:
                L2.append(L1[i])
                i += 1
        L2.extend(L[j:]+L1[i:])
        return L2
    def __init__(self, data):
        self.n = len(data)
        self.tree = [SortedList() for _ in range(2 * self.n)]
        self.build(data)
    def build(self, data):
        for i in range(self.n):
            self.tree[self.n + i].insert(data[i])
        for i in range(self.n - 1, 0, -1):
            self.tree[i] = SortedList(MergeSortTree.merge(self.tree[i << 1],self.tree[i<<1|1]))
    def query(self, left, right,l,r):
        res = 0
        left += self.n
        right += self.n
        while left < right:
            if left & 1:
                res += self.func(self.tree[left], l, r);left += 1
            if right & 1:
                right -= 1;res += self.func(self.tree[right], l, r)
            left >>= 1;right >>= 1
        return res
    def func(self, L, l, r):
        return L.bisect_right(r)-L.bisect_left(l)
    def update(self,pos,value):
        pos += self.n;old_val = self.tree[pos][0]
        self.tree[pos].remove(old_val)
        self.tree[pos].add(value)
        while pos:
            pos >>= 1
            self.tree[pos].remove(old_val)
            self.tree[pos].add(value)
\end{lstlisting}

Persistent Segment Tree
\begin{lstlisting}[language=Python]

class PersistentSegmentTree:
    # Everything is 0 based indexing and [l,r]
    class Node:
        def __init__(self, value=0, left=None, right=None):
            self.value = value
            self.left = left
            self.right = right
    @staticmethod
    def func(a, b):
        return a+b
    def __init__(self, data):
        self.n = len(data)
        self.versions = []
        self.versions.append(self._build(data))
    def _build(self, data):
        stack = [(0, self.n - 1, False)]
        nodes = {}
        while stack:
            left, right, visited = stack.pop()
            if left == right:
                nodes[(left, right)] = self.Node(data[left])
            else:
                if visited:
                    mid = (left + right) // 2
                    left_child = nodes[(left, mid)]
                    right_child = nodes[(mid + 1, right)]
                    nodes[(left, right)] = self.Node(self.func(left_child.value,right_child.value),left_child, right_child)
                else:
                    stack.append((left, right, True))
                    mid = (left + right) // 2
                    stack.append((mid + 1, right, False))
                    stack.append((left, mid, False))
        return nodes[(0, self.n - 1)]
    def update(self, version, pos, value):
        # Returns the root to the new tree,
        # now you must add it to versions yourself
        # seg.versions[version] = seg.update(version, pos, value)
        old_root = self.versions[version]
        stack, path = [(old_root, 0, self.n - 1)], []
        while stack:
            node, left, right = stack.pop()
            path.append((node, left, right))
            if left == right:
                break
            mid = (left + right) // 2
            if pos <= mid:
                stack.append((node.left, left, mid))
            else:
                stack.append((node.right, mid + 1, right))
        new_nodes = {}
        for node, left, right in reversed(path):
            if left == right:
                k = value
                # k = self.func(node.value,value) # if i want to update
                new_nodes[(left, right)] = self.Node(k)
            else:
                mid = (left + right) // 2
                left_child = new_nodes.get((left, mid), node.left)
                right_child = new_nodes.get((mid + 1, right), node.right)
                new_nodes[(left, right)] = self.Node(self.func(left_child.value,right_child.value),left_child, right_child)
        return new_nodes[(0, self.n - 1)]
    def create_version(self, version, pos, value):
        new_root = self.update(version, pos, value)
        self.versions.append(new_root)
        return len(self.versions) - 1
    def query(self, version, ql, qr):
        node, left, right = self.versions[version], 0, self.n - 1
        stack = [(node, left, right)]
        result = 0 # change this depending on the problem
        while stack:
            node, left, right = stack.pop()
            if ql > right or qr < left:
                continue
            if ql <= left and right <= qr:
                result = self.func(result,node.value)
            else:
                mid = (left + right) // 2
                stack.append((node.left, left, mid))
                stack.append((node.right, mid + 1, right))
        return result
\end{lstlisting}

Suffix Automaton
\begin{lstlisting}[language=Python]
class State:
    def __init__(self):
        self.next = {}
        self.link = -1
        self.len = 0
        self.first_pos = -1
        self.occurrence = 0

class SuffixAutomaton:
    def __init__(self, s):
        self.s = s
        self.states = [State()]
        self.size = 1
        self.last = 0
        for ch in s:
            self.add(ch)
        self._prepare_occurrences() # comment out if taking time
        self._count_substrings()
    def add(self, ch):
        p = self.last
        cur = self.size
        self.states.append(State())
        self.size += 1
        self.states[cur].len = self.states[p].len + 1
        self.states[cur].first_pos = self.states[cur].len - 1
        self.states[cur].occurrence = 1
        while p != -1 and ch not in self.states[p].next:
            self.states[p].next[ch] = cur
            p = self.states[p].link
        if p == -1:
            self.states[cur].link = 0
        else:
            q = self.states[p].next[ch]
            if self.states[p].len + 1 == self.states[q].len:
                self.states[cur].link = q
            else:
                clone = self.size
                self.states.append(State())
                self.size += 1
                self.states[clone].len = self.states[p].len + 1
                self.states[clone].next = self.states[q].next.copy()
                self.states[clone].link = self.states[q].link
                self.states[clone].first_pos = self.states[q].first_pos
                while p != -1 and self.states[p].next[ch] == q:
                    self.states[p].next[ch] = clone
                    p = self.states[p].link
                self.states[q].link = self.states[cur].link = clone
        self.last = cur
    def _prepare_occurrences(self):
        order = sorted(range(self.size), key=lambda x: -self.states[x].len)
        for i in order:
            if self.states[i].link != -1:
                self.states[self.states[i].link].occurrence += self.states[i].occurrence
    def _count_substrings(self):
        self.dp = [0] * self.size
        for i in range(self.size):
            self.dp[i] = 1
        order = sorted(range(self.size), key=lambda x: self.states[x].len)
        for u in reversed(order):
            for v in self.states[u].next.values():
                self.dp[u] += self.dp[v]
    def is_substring(self, s):
        current = 0
        for ch in s:
            if ch not in self.states[current].next:
                return False
            current = self.states[current].next[ch]
        return True
    def count_occurrences(self, s):
        current = 0
        for ch in s:
            if ch not in self.states[current].next:
                return 0
            current = self.states[current].next[ch]
        return self.states[current].occurrence
    def count_distinct_substrings(self):
        return sum(self.states[i].len - self.states[self.states[i].link].len for i in range(1, self.size))
    def kth_lex_substring(self, k):
        # kth distinct substring
        result = []
        current = 0
        while k:
            for ch in sorted(self.states[current].next):
                next_state = self.states[current].next[ch]
                if self.dp[next_state] < k:
                    k -= self.dp[next_state]
                else:
                    result.append(ch)
                    k -= 1
                    current = next_state
                    break
        return ''.join(result)
    def enumerate_all_substrings(self):
        result = []
        def dfs(state, path):
            for ch in sorted(self.states[state].next):
                next_state = self.states[state].next[ch]
                result.append(path + ch)
                dfs(next_state, path + ch)
        dfs(0, "")
        return result
    def longest_common_substring(self, t):
        v = 0;l = 0;best = 0;bestpos = 0
        for i in range(len(t)):
            while v and t[i] not in self.states[v].next:
                v = self.states[v].link
                l = self.states[v].len
            if t[i] in self.states[v].next:
                v = self.states[v].next[t[i]]
                l += 1
            if l > best:
                best = l
                bestpos = i
        return t[bestpos - best + 1:bestpos + 1]
    def all_occurrences(self, s):
        current = 0
        for ch in s:
            if ch not in self.states[current].next:
                return []
            current = self.states[current].next[ch]
        positions = []
        def collect(state):
            if self.states[state].occurrence:
                pos = self.states[state].first_pos - len(s) + 1
                positions.append(pos)
            for v in self.states[state].next.values():
                collect(v)
        collect(current)
        return sorted(set(positions))
    def missing_sub(self):
        visited = set()
        q = deque([(0, "")])
        while q:
            state, path = q.popleft()
            for c in map(chr, range(97, 123)):
                if c not in self.states[state].next:
                    return path + c
                next_state = self.states[state].next[c]
                if (next_state, path + c) not in visited:
                    visited.add((next_state, path + c))
                    q.append((next_state, path + c))
        return None
\end{lstlisting}

Segment Tree
\begin{lstlisting}[language=Python]
class SegmentTree:
    @staticmethod
    def func(a, b):
        # Change this function depending upon needs
        return max(a, b)
    def __init__(self, data):
        self.n = len(data)
        self.tree = [0] * (self.n<<1)
        self.build(data)
    def build(self, data):
        for i in range(self.n):
            self.tree[self.n + i] = data[i]
        for i in range(self.n - 1, 0, -1):
            self.tree[i] = self.func(self.tree[i<<1], self.tree[(i<<1) + 1])
    def update(self, pos, value):
        # Update the value at the leaf node
        pos += self.n
        # For updating
        self.tree[pos] = value
        # self.tree[pos] += value
        # If you want to add rather than update
        while pos > 1:
            pos >>= 1
            self.tree[pos] = self.func(self.tree[pos<<1], self.tree[(pos<<1) + 1])
    def query(self, left, right):
        # Query the maximum value in the range [left, right)
        left += self.n
        right += self.n
        # Change the initializer depending upon the self.func
        max_val = float('-inf')
        ##
        while left < right:
            if left&1:
                max_val = self.func(max_val, self.tree[left])
                left += 1
            if right&1:
                right -= 1
                max_val = self.func(max_val, self.tree[right])
            left >>= 1
            right >>= 1
        return max_val
    def __repr__(self):
        values = [str(self.query(i, i + 1)) for i in range(self.n)]
        return f"Seg[{', '.join(values)}]"
\end{lstlisting}

Sparse Table
\begin{lstlisting}[language=Python]
class SparseTable:
    @staticmethod
    def func(a,b):
        # func(a,a) should be a
        return gcd(a,b)
    def __init__(self, arr):
        self.n = len(arr)
        self.table = [[0 for i in range(int((math.log(self.n, 2)+1)))] for j in range(self.n)]
        self.build(arr)
    def build(self, arr):
        for i in range(0, self.n):
            self.table[i][0] = arr[i]
        j = 1
        while (1 << j) <= self.n:
            i = 0
            while i <= self.n - (1 << j):
                self.table[i][j] = self.func(self.table[i][j - 1], self.table[i + (1 << (j - 1))][j - 1])
                i += 1
            j += 1
    def query(self, L, R):
        # query from [L,R]
        j = int(math.log2(R - L + 1))
        return self.func(self.table[L][j], self.table[R - (1 << j) + 1][j])
\end{lstlisting}

BitArray
\begin{lstlisting}[language=Python]
class BitSet:
    ADDRESS_BITS_PER_WORD = 12
    BITS_PER_WORD = 1 << ADDRESS_BITS_PER_WORD
    WORD_MASK = -1
    def __init__(self, sz):
        self.sz = sz
        self.words = [0] * (self._wordIndex(sz - 1) + 1)
    def _wordIndex(self, bitIndex):
        if bitIndex >= self.sz:
            raise ValueError("out of bound index", bitIndex)
        return bitIndex >> BitSet.ADDRESS_BITS_PER_WORD
    def flip(self, bitIndex):
        wordIndex = self._wordIndex(bitIndex)
        self.words[wordIndex] ^= 1 << (bitIndex % BitSet.BITS_PER_WORD)
    def flip_range(self, l, r):
        startWordIndex = self._wordIndex(l)
        endWordIndex = self._wordIndex(r)
        firstWordMask = BitSet.WORD_MASK << (l % BitSet.BITS_PER_WORD)
        rem = (r+1) % BitSet.BITS_PER_WORD
        lastWordMask = BitSet.WORD_MASK if rem == 0 else ~(BitSet.WORD_MASK << rem)
        if startWordIndex == endWordIndex:
            self.words[startWordIndex] ^= (firstWordMask & lastWordMask)
        else:
            self.words[startWordIndex] ^= firstWordMask
            for i in range(startWordIndex + 1, endWordIndex):
                self.words[i] ^= BitSet.WORD_MASK
            self.words[endWordIndex] ^= lastWordMask
    def __setitem__(self, bitIndex, value):
        wordIndex = self._wordIndex(bitIndex)
        if value:
            self.words[wordIndex] |= 1 << (bitIndex % BitSet.BITS_PER_WORD)
        else:
            self.words[wordIndex] &= ~(1 << (bitIndex % BitSet.BITS_PER_WORD))
    def __getitem__(self, bitIndex):
        wordIndex = self._wordIndex(bitIndex)
        return self.words[wordIndex] & (1 << (bitIndex % BitSet.BITS_PER_WORD)) != 0
    def nextSetBit(self, fromIndex):
        wordIndex = self._wordIndex(fromIndex)
        word = self.words[wordIndex] & (BitSet.WORD_MASK << (fromIndex % BitSet.BITS_PER_WORD))

        while True:
            if word != 0:
                return wordIndex * BitSet.BITS_PER_WORD + (word & -word).bit_length() - 1
            wordIndex += 1
            if wordIndex > len(self.words) - 1:
                return -1
            word = self.words[wordIndex]
    def nextClearBit(self, fromIndex):
        wordIndex = self._wordIndex(fromIndex)
        word = ~self.words[wordIndex] & (BitSet.WORD_MASK << (fromIndex % BitSet.BITS_PER_WORD))

        while True:
            if word != 0:
                index = wordIndex * BitSet.BITS_PER_WORD + (word & -word).bit_length() - 1
                return index if index < self.sz else - 1
            wordIndex += 1
            if wordIndex > len(self.words) - 1:
                return -1
            word = ~self.words[wordIndex]
    def lastSetBit(self):
        wordIndex = len(self.words) - 1
        word = self.words[wordIndex]

        while wordIndex >= 0:
            if word != 0:
                return wordIndex * BitSet.BITS_PER_WORD + (word.bit_length() - 1 if word > 0 else  BitSet.BITS_PER_WORD - 1)
            wordIndex -= 1
            word = self.words[wordIndex]
        return -1
    def __str__(self):
        res = []
        st = 0
        while True:
            i = self.nextSetBit(st)
            if i != -1:
                res += [0] * (i - st)
                j = self.nextClearBit(i)
                if j != -1:
                    res += [1] * (j-i)
                    st = j
                else:
                    res += [1] * (self.sz - i)
                    break
            else:
                res += [0] * (self.sz - st)
                break

        return "".join(str(v) for v in res)
    def __repr__(self):
        return "Bitset(%s)" % str(self)
    def __iter__(self):
        for i in self[:]:
            yield i
    def __len__(self):
        return self.sz
    def __or__(self, other):
        if self.sz != other.sz:
            raise ValueError("BitSets must be of equal size")
        res = BitSet(self.sz)
        res.words = [a | b for a, b in zip(self.words, other.words)]
        return res
    def __add__(self, other):
        if self.sz != other.sz:
            raise ValueError("BitSets must be of equal size")
        res = BitSet(self.sz)
        carry = 0
        for i in range(len(self.words)):
            total = self.words[i] + other.words[i] + carry
            res.words[i] = total & BitSet.WORD_MASK
            carry = total >> BitSet.BITS_PER_WORD
        return res
    def __and__(self, other):
        if self.sz != other.sz:
            raise ValueError("BitSets must be of equal size")
        res = BitSet(self.sz)
        res.words = [a & b for a, b in zip(self.words, other.words)]
        return res
    def __xor__(self, other):
        if self.sz != other.sz:
            raise ValueError("BitSets must be of equal size")
        res = BitSet(self.sz)
        res.words = [a ^ b for a, b in zip(self.words, other.words)]
        return res
    def __invert__(self):
        res = BitSet(self.sz)
        res.words = [~a & BitSet.WORD_MASK for a in self.words]
        return res
    def add(self, val):
        self.flip_range(val, self.nextClearBit(val))
    def rem(self, val):
        self.flip_range(val, self.nextSetBit(val))
\end{lstlisting}

\section{Graphs}

Basics
\begin{lstlisting}[language=Python]
def dfs(d,start=0):
    n = len(d)
    visited = [False]*n
    dp = [0]*n
    finished = [False]*n
    stack = [start]
    while stack:
        start = stack[-1]
        if not visited[start]:
            visited[start] = True
            for child in d[start]:
                if not visited[child]:
                    stack.append(child)
        else:
            stack.pop()
            dp[start] += 1
            for child in d[start]:
                if finished[child]:
                    dp[start] += dp[child]
            finished[start] = True
    return dp

def dijkstra(d,start=0):
    n = len(d)
    dist = [float("inf")]*n
    # parents = [-1]*n
    dist[start] = 0
    queue = [(0, start)]
    while queue:
        path_len, v = heappop(queue)
        if path_len == dist[v]:
            for w, edge_len in d[v]:
                new_dist = edge_len+path_len
                if new_dist<dist[w]:
                    dist[w] = new_dist
                    # parents[w] = v
                    heappush(queue, (new_dist, w))
    return dist
\end{lstlisting}

Others
\begin{lstlisting}[language=Python]
def euler_path(d):
    start = [1]
    ans = []
    while start:
        cur = start[-1]
        if len(d[cur])==0:
            ans.append(start.pop())
            continue
        k1 = d[cur].pop()
        d[k1].remove(cur) # if undirected
        start.append(k1)
    return ans

def floyd_warshall(n, edges):
    dist = [[0 if i == j else float("inf") for i in range(n)] for j in range(n)]
    pred = [[None] * n for _ in range(n)]
    for u, v, d in edges:
        dist[u][v] = d
        pred[u][v] = u
    for k in range(n):
        for i in range(n):
            for j in range(n):
                if dist[i][k] + dist[k][j] < dist[i][j]:
                    dist[i][j] = dist[i][k] + dist[k][j]
                    pred[i][j] = pred[k][j]
    return dist, pred

def toposort(graph):
    res, found = [], [0] * len(graph)
    stack = list(range(len(graph)))
    while stack:
        node = stack.pop()
        if node < 0:
            res.append(~node)
        elif not found[node]:
            found[node] = 1
            stack.append(~node)
            stack += graph[node]
    for node in res:
        if any(found[nei] for nei in graph[node]):
            return None
        found[node] = 0
    return res[::-1]
\end{lstlisting}

Advance
\begin{lstlisting}[language=Python]
class AuxiliaryTree:
    def __init__(self, edge, root = 0):
        self.n = len(edge)
        self.order = [-1] * self.n
        self.path = [-1] * (self.n-1)
        self.depth = [0] * self.n
        if self.n == 1: return
        parent = [-1] * self.n
        que = [root]
        t = -1
        while que:
            u = que.pop()
            self.path[t] = parent[u]
            t += 1
            self.order[u] = t
            for v in edge[u]:
                if self.order[v] == -1:
                    que.append(v)
                    parent[v] = u
                    self.depth[v] = self.depth[u] + 1
        self.n -= 1
        self.h = self.n.bit_length()
        self.data = [0] * (self.n * self.h)
        self.data[:self.n] = [self.order[u] for u in self.path]
        for i in range(1, self.h):
            for j in range(self.n - (1<<i) + 1):
                self.data[i*self.n + j] = min(self.data[(i-1)*self.n + j], self.data[(i-1)*self.n + j+(1<<(i-1))])

    def lca(self, u, v):
        if u == v: return u
        l = self.order[u]
        r = self.order[v]
        if l > r:
            l,r = r,l
        level = (r - l).bit_length() - 1
        return self.path[min(self.data[level*self.n + l], self.data[level*self.n + r-(1<<level)])]

    def dis(self, u, v):
        if u == v: return 0
        l = self.order[u]
        r = self.order[v]
        if l > r:
            l,r = r,l
        level = (r - l).bit_length() - 1
        p = self.path[min(self.data[level*self.n + l], self.data[level*self.n + r-(1<<level)])]
        return self.depth[u] + self.depth[v] - 2 * self.depth[p]

    def make(self, vs):
        k = len(vs)
        vs.sort(key = self.order.__getitem__)

        par = dict()
        edge = dict()
        edge[vs[0]] = []

        st = [vs[0]]

        for i in range(k - 1):
            l = self.order[vs[i]]
            r = self.order[vs[i+1]]
            level = (r - l).bit_length() - 1
            w = self.path[min(self.data[level*self.n + l], self.data[level*self.n + r-(1<<level)])]
            if w != vs[i]:
                p = st.pop()
                while st and self.depth[w] < self.depth[st[-1]]:
                    par[p] = st[-1]
                    edge[st[-1]].append(p)
                    p = st.pop()

                if not st or st[-1] != w:
                    st.append(w)
                    edge[w] = [p]
                else:
                    edge[w].append(p)
                par[p] = w

            st.append(vs[i+1])
            edge[vs[i+1]] = []

        for i in range(len(st) - 1):
            edge[st[i]].append(st[i+1])
            par[st[i+1]] = st[i]

        par[st[0]] = -1
        return st[0], edge, par

class binary_lift:
    def __init__(self, graph, f=max, root=0, flag=False):
        n = len(graph)
        parent = [-1] * (n + 1)
        depth = self.depth = [-1] * n
        bfs = [root]
        depth[root] = 0
        data = [0]*n
        for node in bfs:
            # for nei,w in graph[node]:
            for nei in graph[node]:
                if depth[nei] == -1:
                    # data[nei] = w
                    parent[nei] = node
                    depth[nei] = depth[node] + 1
                    bfs.append(nei)
        parent = self.parent = [parent]
        self.f = f
        if flag:
            data = self.data = [data]
            for _ in range(max(depth).bit_length()):
                old_data = data[-1]
                old_parent = parent[-1]
                data.append([f(val, old_data[p]) for val,p in zip(old_data, old_parent)])
                parent.append([old_parent[p] for p in old_parent])
        else:
            for _ in range(max(depth).bit_length()):
                old_parent = parent[-1]
                parent.append([old_parent[p] for p in old_parent])
    def lca(self, a, b):
        depth = self.depth
        parent = self.parent
        if depth[a] < depth[b]:
            a,b = b,a
        d = depth[a] - depth[b]
        for i in range(d.bit_length()):
            if (d >> i) & 1:
                a = parent[i][a]
        for i in range(depth[a].bit_length())[::-1]:
            if parent[i][a] != parent[i][b]:
                a = parent[i][a]
                b = parent[i][b]
        if a != b:
            return parent[0][a]
        else:
            return a
    def distance(self, a, b):
        return self.depth[a] + self.depth[b] - 2 * self.depth[self.lca(a,b)]
    def kth_ancestor(self, a, k):
        parent = self.parent
        if self.depth[a] < k:
            return -1
        for i in range(k.bit_length()):
            if (k >> i) & 1:
                a = parent[i][a]
        return a
    def __call__(self, a, b, c=0):
        depth = self.depth
        parent = self.parent
        data = self.data
        f = self.f
        c = self.lca(a, b)
        val = c
        for x,d in (a, depth[a] - depth[c]), (b, depth[b] - depth[c]):
            for i in range(d.bit_length()):
                if (d >> i) & 1:
                    val = f(val, data[i][x])
                    x = parent[i][x]
        return val
\end{lstlisting}

Graph-Flatten
\begin{lstlisting}[language=Python]
def dfs(graph):
    starttime = [[0,0] for i in range(len(graph))]
    time = 0
    stack = [(0,-1,0)]
    while stack:
        cur, prev, state = stack.pop()
        if state == 0:
            starttime[cur][0] = time
            time += 1
            stack.append((cur, prev, 1))
            for neighbor in graph[cur]:
                if neighbor == prev:
                    continue
                stack.append((neighbor, cur, 0))
        elif state == 1:
            starttime[cur][1] = time
    return starttime
\end{lstlisting}

2-sat
\begin{lstlisting}[language=Python]
def find_SCC(graph):
    SCC, S, P = [], [], []
    depth = [0] * len(graph)
    stack = list(range(len(graph)))
    while stack:
        node = stack.pop()
        if node < 0:
            d = depth[~node] - 1
            if P[-1] > d:
                SCC.append(S[d:])
                del S[d:], P[-1]
                for node in SCC[-1]:
                    depth[node] = -1
        elif depth[node] > 0:
            while P[-1] > depth[node]:
                P.pop()
        elif depth[node] == 0:
            S.append(node)
            P.append(len(S))
            depth[node] = len(S)
            stack.append(~node)
            stack += graph[node]
    return SCC[::-1]

class TwoSat:
    def __init__(self, n):
        self.n = n
        self.graph = [[] for _ in range(2 * n)]
    def negate(self, x):
        return x+self.n if x<self.n else x-self.n
    def _imply(self, x, y):
        # agar x hoga , toh y hoga
        self.graph[x].append(y)
        self.graph[self.negate(y)].append(self.negate(x))
    def either(self, x, y):
        # koi ek true ho sakta hain ya dono bhi
        self._imply(self.negate(x),y)
        self._imply(self.negate(y),x)
    def set(self, x):
        self._imply(self.negate(x),x)
    def solve(self):
        SCC = find_SCC(self.graph)
        order = [0] * (2 * self.n)
        for i, comp in enumerate(SCC):
            for x in comp:
                order[x] = i
        for i in range(self.n):
            if order[i] == order[self.negate(i)]:
                return False, None
        return True, [+(order[i] > order[self.negate(i)]) for i in range(self.n)]
\end{lstlisting}

Bridges and Aritculation Points
\begin{lstlisting}[language=Python]
    def find_bridges(adj):
    # returns all bridges
    bridges = []
    n = len(adj)
    timer = 0
    visited = [False]*n
    tin = [-1]*n
    low = [-1]*n
    for start in range(n):
        if visited[start]:
            continue
        stack = [(start, -1, 0, False)]
        visited[start] = True
        tin[start] = low[start] = timer
        timer += 1
        while stack:
            v, parent, idx, backtrack = stack.pop()
            if backtrack:
                to = adj[v][idx]
                low[v] = min(low[v], low[to])
                if low[to] > tin[v]:
                    bridges.append((v,to))
                continue
            if idx < len(adj[v]):
                to = adj[v][idx]
                stack.append((v, parent, idx + 1, False))
                if to == parent:
                    continue
                if visited[to]:
                    low[v] = min(low[v], tin[to])
                else:
                    visited[to] = True
                    tin[to] = low[to] = timer
                    timer += 1
                    stack.append((v, parent, idx, True))
                    stack.append((to, v, 0, False))
    return bridges

def bridges_on_path(adj):
    # returns all bridges on path from 1 to n
    n = len(adj)
    timer = 0
    visited = [False]*n
    tin = [-1]*n
    low = [-1]*n
    bridges = []
    for start in range(n):
        if visited[start]:
            continue
        stack = [(start, -1, 0, False)]
        visited[start] = True
        tin[start] = low[start] = timer
        timer += 1
        while stack:
            v, parent, idx, backtrack = stack.pop()
            if backtrack:
                to = adj[v][idx]
                low[v] = min(low[v], low[to])
                if low[to] > tin[v]:
                    bridges.append((v, to))
                continue
            if idx < len(adj[v]):
                to = adj[v][idx]
                stack.append((v, parent, idx+1, False))
                if to == parent:
                    continue
                if visited[to]:
                    low[v] = min(low[v], tin[to])
                else:
                    visited[to] = True
                    tin[to] = low[to] = timer
                    timer += 1
                    stack.append((v, parent, idx, True))
                    stack.append((to, v, 0, False))
    comp_id = [-1]*n
    comp = 0
    bridge_set = set(bridges)
    for i in range(n):
        if comp_id[i] != -1:
            continue
        stack = [i]
        comp_id[i] = comp
        while stack:
            v = stack.pop()
            for to in adj[v]:
                if comp_id[to] == -1 and (v, to) not in bridge_set and (to, v) not in bridge_set:
                    comp_id[to] = comp
                    stack.append(to)
        comp += 1

    tree = [[] for _ in range(comp)]
    for u, v in bridges:
        cu, cv = comp_id[u], comp_id[v]
        tree[cu].append((cv, (u, v)))
        tree[cv].append((cu, (u, v)))

    c1, cN = comp_id[0], comp_id[n-1]
    if c1 == cN:
        return []
    parent = {c1: None}
    edge_used = {c1: None}
    stack = [c1]
    while stack:
        u = stack.pop()
        for v, e in tree[u]:
            if v not in parent:
                parent[v] = u
                edge_used[v] = e
                stack.append(v)
    path_bridges = []
    cur = cN
    while cur != c1:
        path_bridges.append(edge_used[cur])
        cur = parent[cur]
    return [(u,v) for u, v in path_bridges]

def lowlink(edge):
    n = len(edge)
    parent = [-1] * n
    visited = [False] * n
    for s in range(n):
        if not visited[s]:
            que = [s]
            while que:
                now = que.pop()
                if visited[now]: continue
                visited[now] = True
                for nxt in edge[now]:
                    if not visited[nxt]:
                        parent[nxt] = now
                        que.append(nxt)
    order = [-1] * n
    low = [-1] * n
    is_articulation = [False] * n
    articulation = []
    bridge = []
    def dfs(s):
        idx = 0
        cnt = 0
        que = [~s,s]
        while que:
            now = que.pop()
            if now >= 0:
                order[now] = low[now] = idx
                idx += 1
                for nxt in edge[now]:
                    if parent[nxt] == now:
                        que.append(~nxt)
                        que.append(nxt)
                    elif parent[now] != nxt and order[nxt] != -1:
                        low[now] = min(low[now], order[nxt])
            else:
                now = ~now
                par = parent[now]
                if par == s: cnt += 1
                if now == s:
                    is_articulation[now] |= (cnt >= 2)
                    if is_articulation[now]:
                        articulation.append(now)
                    return
                if is_articulation[now]:
                    articulation.append(now)
                if now != parent[par]:
                    low[par] = min(low[par], low[now])
                is_articulation[par] |= (par != s) and (order[par] <= low[now])
                if order[par] < low[now]:
                    bridge.append((par, now))
    for i in range(n):
        if parent[i] == -1:
            dfs(i)
    return articulation, bridge

def find_2ecc(edges,d):
    # returns a new graph, in which two nodes are connected
    # if and only if they are part of same cycle.
    _,bridges = lowlink(d)
    newd = [[] for i in range(len(d))]
    bridges = set((w(i[0]),w(i[1])) for i in bridges)
    for u,v in edges:
        if (w(u),w(v)) not in bridges and (w(v),w(u)) not in bridges:
            newd[u].append(v)
            newd[v].append(u)
    return newd
\end{lstlisting}


\section{Strings}

String Hashing
\begin{lstlisting}[language=Python]
HMOD = 2147483647
HBASE1 = random.randrange(HMOD)
HBASE2 = random.randrange(HMOD)
class Hashing:
    def __init__(self, s, mod=HMOD, base1=HBASE1, base2=HBASE2):
        self.mod, self.base1, self.base2 = mod, base1, base2
        self._len = _len = len(s)
        f_hash, f_pow = [0] * (_len + 1), [1] * (_len + 1)
        for i in range(_len):
            f_hash[i + 1] = (base1 * f_hash[i] + s[i]) % mod
            f_pow[i + 1] = base1 * f_pow[i] % mod
        self.f_hash, self.f_pow = f_hash, f_pow
    def hashed(self, start, stop):
        return (self.f_hash[stop] - self.f_pow[stop - start] * self.f_hash[start]) % self.mod

\end{lstlisting}

String Functions
\begin{lstlisting}[language=Python]
def LCPArray(L,SA=None):
    # Longest Common prefix in between S[i:] and S[i+1:]
    if not SA:
        SA = SuffArr(L)
    n = len(L)
    rank = [0] * n
    for i in range(n):
        rank[SA[i]] = i
    LCP = [0] * (n - 1)
    k = 0
    for i in range(n):
        l = rank[i]
        if l==n-1:
            k = 0
            continue
        j = SA[l+1]
        while i+k<n and L[i+k] == L[j+k]:
            k += 1
        LCP[l] = k
        k -= k > 0
    return LCP

def z_function(S):
    # return: the Z array, where Z[i] = length of the longest common prefix of S[i:] and S
    n = len(S)
    Z = [0] * n
    l = r = 0
    for i in range(1, n):
        z = Z[i - l]
        if i + z >= r:
            z = max(r - i, 0)
            while i + z < n and S[z] == S[i + z]:
                z += 1
            l, r = i, i + z
        Z[i] = z
    Z[0] = n
    return Z

def manacher(s):
    # returns longest palindrome in s
    t = '#' + '#'.join(s) + '#'
    n = len(t)
    L = [0] * n
    c = r = 0
    ml = 0
    mc = 0
    for i in range(n):
        mirror = 2*c-i
        if i<r:
            L[i] = min(r-i,L[mirror])
        a = i+L[i]+1
        b = i-L[i]-1
        while a<n and b>=0 and t[a]==t[b]:
            L[i] += 1
            a += 1
            b -= 1
        if i+L[i]>r:
            c = i
            r = i+L[i]
        if L[i]>ml:
            ml = L[i]
            mc = i
    start = (mc-ml)//2
    return s[start:start+ml]

def SuffArray(s):
    # Starting position of ith suffix
    # in lexographic order
    s += "$"
    n = len(s)
    k = 0
    rank = [ord(c) for c in s]
    tmp = [0] * n
    sa = list(range(n))
    def sort_key(i):
        return (rank[i], rank[i + (1 << k)] if i + (1 << k) < n else -1)
    while True:
        sa.sort(key=sort_key)
        tmp[sa[0]] = 0
        for i in range(1, n):
            tmp[sa[i]] = tmp[sa[i-1]] + (sort_key(sa[i-1]) != sort_key(sa[i]))
        rank, tmp = tmp, rank
        k += 1
        if (1 << k) >= n:
            break
    return sa
\end{lstlisting}


\section{FFT}
FFT
\begin{lstlisting}[language=Python]
class FastFFT:
    # This is a bit faster, one log n factor is less but it's accuracy is not 100%
    # use this when coeffecient does not matter(set coeffecient to 1 repeatedly)
    # or you could just risk it :)
    def __init__(self, MOD=998244353):
        FastFFT.MOD = MOD
        # g = self.primitive_root_constexpr()
        g = 3
        ig = pow(g, FastFFT.MOD - 2, FastFFT.MOD)
        FastFFT.W = [pow(g, (FastFFT.MOD - 1) >> i, FastFFT.MOD) for i in range(30)]
        FastFFT.iW = [pow(ig, (FastFFT.MOD - 1) >> i, FastFFT.MOD) for i in range(30)]
    def primitive_root_constexpr(self):
        if FastFFT.MOD == 998244353:
            return 3
        elif FastFFT.MOD == 200003:
            return 2
        elif FastFFT.MOD == 167772161:
            return 3
        elif FastFFT.MOD == 469762049:
            return 3
        elif FastFFT.MOD == 754974721:
            return 11
        divs = [0] * 20
        divs[0] = 2
        cnt = 1
        x = (FastFFT.MOD - 1) // 2
        while x % 2 == 0:
            x //= 2
        i = 3
        while i * i <= x:
            if x % i == 0:
                divs[cnt] = i
                cnt += 1
                while x % i == 0:
                    x //= i
            i += 2
        if x > 1:
            divs[cnt] = x
            cnt += 1
        g = 2
        while 1:
            ok = True
            for i in range(cnt):
                if pow(g, (FastFFT.MOD - 1) // divs[i], FastFFT.MOD) == 1:
                    ok = False
                    break
            if ok:
                return g
            g += 1
    def fft(self, k, f):
        for l in range(k, 0, -1):
            d = 1 << l - 1
            U = [1]
            for i in range(d):
                U.append(U[-1] * FastFFT.W[l] % FastFFT.MOD)
            for i in range(1 << k - l):
                for j in range(d):
                    s = i * 2 * d + j
                    f[s], f[s + d] = (f[s] + f[s + d]) % FastFFT.MOD, U[j] * (f[s] - f[s + d]) % FastFFT.MOD
    def ifft(self, k, f):
        for l in range(1, k + 1):
            d = 1 << l - 1
            for i in range(1 << k - l):
                u = 1
                for j in range(i * 2 * d, (i * 2 + 1) * d):
                    f[j+d] *= u
                    f[j], f[j + d] = (f[j] + f[j + d]) % FastFFT.MOD, (f[j] - f[j + d]) % FastFFT.MOD
                    u = u * FastFFT.iW[l] % FastFFT.MOD
    def convolve(self, A, B):
        n0 = len(A) + len(B) - 1
        k = (n0).bit_length()
        n = 1 << k
        A += [0] * (n - len(A))
        B += [0] * (n - len(B))
        self.fft(k, A)
        self.fft(k, B)
        A = [a * b % FastFFT.MOD for a, b in zip(A, B)]
        self.ifft(k, A)
        inv = pow(n, FastFFT.MOD - 2, FastFFT.MOD)
        del A[n0:]
        for i in range(n0):
            A[i] = (A[i]*inv)%FastFFT.MOD
        return A

class FFT:
    def __init__(self, MOD=998244353,MOD1=469762049):
        FFT.MOD = MOD
        FFT.MOD1 = MOD1
        FFT.MOD2 = pow(MOD,MOD1-2,MOD1)
        FFT.mod_inv = (self.XT_GCD(MOD,MOD1)[1])%MOD1
        # g = self.primitive_root_constexpr()
        g = 3
        ig = pow(g, FFT.MOD - 2, FFT.MOD)
        ig1 = pow(g, FFT.MOD1 - 2, FFT.MOD1)
        FFT.W = [pow(g, (FFT.MOD - 1) >> i, FFT.MOD) for i in range(30)]
        FFT.W1 = [pow(g, (FFT.MOD1 - 1) >> i, FFT.MOD1) for i in range(30)]
        FFT.iW = [pow(ig, (FFT.MOD - 1) >> i, FFT.MOD) for i in range(30)]
        FFT.iW1 = [pow(ig1, (FFT.MOD1 - 1) >> i, FFT.MOD1) for i in range(30)]
    def primitive_root_constexpr(self):
        if FFT.MOD == 998244353:
            return 3
        elif FFT.MOD == 200003:
            return 2
        elif FFT.MOD == 167772161:
            return 3
        elif FFT.MOD == 469762049:
            return 3
        elif FFT.MOD == 754974721:
            return 11
        divs = [0] * 20
        divs[0] = 2
        cnt = 1
        x = (FFT.MOD - 1) // 2
        while x % 2 == 0:
            x //= 2
        i = 3
        while i * i <= x:
            if x % i == 0:
                divs[cnt] = i
                cnt += 1
                while x % i == 0:
                    x //= i
            i += 2
        if x > 1:
            divs[cnt] = x
            cnt += 1
        g = 2
        while 1:
            ok = True
            for i in range(cnt):
                if pow(g, (FFT.MOD - 1) // divs[i], FFT.MOD) == 1:
                    ok = False
                    break
            if ok:
                return g
            g += 1
    def fft(self, k, f,f1):
        for l in range(k, 0, -1):
            d = 1 << l - 1
            U = [(1,1)]
            for i in range(d):
                U.append((U[-1][0] * FFT.W[l] % FFT.MOD,U[-1][1] * FFT.W1[l] % FFT.MOD1))
            for i in range(1 << k - l):
                for j in range(d):
                    s = i * 2 * d + j
                    f[s], f[s + d] = (f[s] + f[s + d]) % FFT.MOD, U[j][0] * (f[s] - f[s + d]) % FFT.MOD
                    f1[s], f1[s + d] = (f1[s] + f1[s + d]) % FFT.MOD1, U[j][1] * (f1[s] - f1[s + d]) % FFT.MOD1
    def ifft(self, k, f,f1):
        for l in range(1, k + 1):
            d = 1 << l - 1
            for i in range(1 << k - l):
                u = 1
                u1 = 1
                for j in range(i * 2 * d, (i * 2 + 1) * d):
                    f[j+d] *= u
                    f[j], f[j + d] = (f[j] + f[j + d]) % FFT.MOD, (f[j] - f[j + d]) % FFT.MOD
                    u = u * FFT.iW[l] % FFT.MOD
                    f1[j+d] *= u1
                    f1[j], f1[j + d] = (f1[j] + f1[j + d]) % FFT.MOD1, (f1[j] - f1[j + d]) % FFT.MOD1
                    u1 = u1 * FFT.iW1[l] % FFT.MOD1
    def XT_GCD(self,a,b):
        if b == 0:
            return a,1,0
        g,x1,y1 = self.XT_GCD(b,a%b)
        x = y1
        y = x1-(a//b)*y1
        return g,x,y
    def CRT(self,a, mod1, b, mod2):
        k = (a+(b-a)*self.mod_inv%mod2*mod1)%(mod1*mod2)
        return k
    def convolve(self, A, B):
        n0 = len(A) + len(B) - 1
        k = (n0).bit_length()
        n = 1 << k
        A += [0] * (n - len(A))
        B += [0] * (n - len(B))
        A1 = A[:]
        B1 = B[:]
        self.fft(k, A,A1)
        self.fft(k, B,B1)
        A = [a * b % FFT.MOD for a, b in zip(A, B)]
        A1 = [a * b % FFT.MOD1 for a, b in zip(A1, B1)]
        self.ifft(k, A,A1)
        inv = pow(n, FFT.MOD - 2, FFT.MOD)
        inv1 = pow(n, FFT.MOD1 - 2, FFT.MOD1)
        del A[n0:]
        for i in range(n0):
            A[i] = self.CRT(A[i]*inv,FFT.MOD,A1[i]*inv1,FFT.MOD1)
        return A
\end{lstlisting}

FWHT
\begin{lstlisting}[language=Python]
def fwht(a, invert=False):
    # In-place Fast Walsh Hadamard Transform for XOR convolution.
    # invert: if True, computes inverse transform
    n = len(a);step = 1
    while step<n:
        for i in range(0,n,step*2):
            for j in range(step):
                a[i+j],a[i+j+step] = a[i+j]+a[i+j+step],a[i+j]-a[i+j+step]
        step *= 2
    if invert:
        for i in range(n):
            a[i] //= n

def fwt_or(a, invert=False):
    n = len(a);step = 1
    while step<n:
        for i in range(0,n,step*2):
            for j in range(step):
                u = a[i+j]
                v = a[i+j+step]
                if not invert:
                    a[i+j+step] = u+v
                else:
                    a[i+j+step] = v-u
        step *= 2

def fwt_and(a, invert=False):
    n = len(a);step = 1
    while step<n:
        for i in range(0,n,step*2):
            for j in range(step):
                u = a[i+j]
                v = a[i+j+step]
                if not invert:
                    a[i+j] = u+v
                else:
                    a[i+j] = u-v
        step *= 2

def convolution(A, B, fwt=fwht):
    # Computes XOR convolution of arrays A and B using FWHT
    # res[i] = summation of A[j]*B[k] such that j xor k = i
    n = 1;x = max(len(A),len(B))
    while n<x:n <<= 1
    n <<= 1
    fa = A + [0] * (n - len(A))
    fb = B + [0] * (n - len(B))
    fwt(fa)
    fwt(fb)
    for i in range(n):
        fa[i] *= fb[i]
    fwt(fa, invert=True)
    return fa
\end{lstlisting}


\section{Miscellaneous}
Bootstrap
\begin{lstlisting}[language=Python]
from types import GeneratorType
def bootstrap(f, stack=[]):
    def wrappedfunc(*args, **kwargs):
        if stack:
            return f(*args, **kwargs)
        to = f(*args, **kwargs)
        while True:
            if type(to) is GeneratorType:
                stack.append(to)
                to = next(to)
            else:
                stack.pop()
                if not stack:
                    break
                to = stack[-1].send(to)
        return to
    return wrappedfunc
\end{lstlisting}

MCMF
\begin{lstlisting}[language=Python]

class MCMF:
    # better when the flow is not constant
    # O(FEV)
    def __init__(self, n):
        self.n = n
        self.g = [[] for _ in range(n)]
    def add_edge(self, u, v, cap, cost):
        self.g[u].append([v, cap, cost, len(self.g[v]), cap])
        self.g[v].append([u, 0, -cost, len(self.g[u]) - 1, 0])
    def min_cost(self, s, t, maxf=float('inf')):
        n = self.n
        flow, cost = 0, 0
        pot = [0] * n
        while flow < maxf:
            dist = [float('inf')] * n
            parent = [(-1, -1)] * n
            dist[s] = 0
            pq = [(0, s)]
            while pq:
                d, u = heappop(pq)
                if d > dist[u]:
                    continue
                for i, (v, cap, w, rev, _) in enumerate(self.g[u]):
                    if cap > 0:
                        nd = d + w + pot[u] - pot[v]
                        if nd < dist[v]:
                            dist[v] = nd
                            parent[v] = (u, i)
                            heappush(pq, (nd, v))
            if dist[t] == float('inf'):
                break
            for i in range(n):
                if dist[i] < float('inf'):
                    pot[i] += dist[i]
            f = maxf - flow
            cur = t
            while cur != s:
                u, i = parent[cur]
                f = min(f, self.g[u][i][1])
                cur = u
            flow += f
            cost += f * pot[t]
            cur = t
            while cur != s:
                u, i = parent[cur]
                v, cap, w, rev, orig = self.g[u][i]
                self.g[u][i][1] -= f
                self.g[cur][rev][1] += f
                cur = u
        return flow, cost
    def used_edges(self):
        res = []
        for u in range(self.n):
            for v, cap, cost, rev, orig in self.g[u]:
                if orig > 0 and cap == 0:
                    res.append((u, v, cost))
        return res

def hungarian(cost):
    # better when the flow is 1 for every edge
    # for i in range(n):print(i+1,match[i]+1)
    # this is the final edges taken
    # O(n3)
    n = len(cost)
    u = [0]*(n+1)
    v = [0]*(n+1)
    p = [0]*(n+1)
    way = [0]*(n+1)
    for i in range(1, n+1):
        p[0] = i
        j0 = 0
        minv = [float('inf')]*(n+1)
        used = [False]*(n+1)
        while True:
            used[j0] = True
            i0 = p[j0]
            delta = float('inf')
            j1 = 0
            for j in range(1, n+1):
                if not used[j]:
                    cur = cost[i0-1][j-1] - u[i0] - v[j]
                    if cur < minv[j]:
                        minv[j] = cur
                        way[j] = j0
                    if minv[j] < delta:
                        delta = minv[j]
                        j1 = j
            for j in range(n+1):
                if used[j]:
                    u[p[j]] += delta
                    v[j] -= delta
                else:
                    minv[j] -= delta
            j0 = j1
            if p[j0] == 0:
                break
        while True:
            j1 = way[j0]
            p[j0] = p[j1]
            j0 = j1
            if j0 == 0:
                break
    match = [-1]*n
    for j in range(1, n+1):
        if p[j] != 0:
            match[p[j]-1] = j-1
    return -v[0], match
\end{lstlisting}

Treap
\begin{lstlisting}[language=Python]

class Node:
    __slots__ = ("val", "priority", "left", "right", "size")
    def __init__(self, val):
        self.val = val
        self.priority = random.random()
        self.left = None
        self.right = None
        self.size = 1
class Treap:
    def __init__(self, iterable=None):
        self.root = None
        if iterable:
            for x in iterable:
                self.root = Treap._merge(self.root, Node(x))
    @staticmethod
    def _size(node):
        return node.size if node else 0
    @staticmethod
    def _update(node):
        if node:
            node.size = 1 + Treap._size(node.left) + Treap._size(node.right)
    @staticmethod
    def _merge(left, right):
        if not left or not right:
            return left or right

        if left.priority > right.priority:
            left.right = Treap._merge(left.right, right)
            Treap._update(left)
            return left
        else:
            right.left = Treap._merge(left, right.left)
            Treap._update(right)
            return right
    def _kth(node, k):
        if not node:
            return None

        left_size = Treap._size(node.left)

        if k < left_size:
            return Treap._kth(node.left, k)
        elif k == left_size:
            return node.val
        else:
            return Treap._kth(node.right, k - left_size - 1)
    def kth(self, k):
        if k < 0 or k >= Treap._size(self.root):
            return None
        return Treap._kth(self.root, k)
    @staticmethod
    def merge(t1, t2):
        t = Treap()
        t.root = Treap._merge(t1.root, t2.root)
        return t
    @staticmethod
    def _split(node, k):
        if not node:
            return (None, None)
        left_size = Treap._size(node.left)
        if k <= left_size:
            left, node.left = Treap._split(node.left, k)
            Treap._update(node)
            return (left, node)
        else:
            node.right, right = Treap._split(node.right, k - left_size - 1)
            Treap._update(node)
            return (node, right)
    def insert(self, pos, val):
        new_node = Node(val)
        left, right = Treap._split(self.root, pos)
        merged = Treap._merge(left, new_node)
        self.root = Treap._merge(merged, right)
    def split(self, k):
        left_root, right_root = Treap._split(self.root, k)
        t1 = Treap()
        t2 = Treap()
        t1.root = left_root
        t2.root = right_root
        return t1, t2
    def __iter__(self):
        stack = []
        node = self.root
        while stack or node:
            while node:
                stack.append(node)
                node = node.left

            node = stack.pop()
            yield node.val
            node = node.right
\end{lstlisting}

DP Optimizations
\begin{lstlisting}[language=Python]
def Knuths(n,C = lambda x: x):
    # conditions
    # For all a<=b<=c<=d:
    #	 cost(a,c)+cost(b,d)<=cost(a,d)+cost(b,c)
    # opt[l][r-1] <= opt[l][r] <= opt[l+1][r]
    dp = [[0]*n for _ in range(n)]
    opt = [[0]*n for _ in range(n)]
    for i in range(n):
        dp[i][i] = 0
        opt[i][i] = i
    for i in range(n-2,-1,-1):
        for j in range(i+1,n):
            cost = C(i, j)
            mini = float('inf')
            L = opt[i][j-1]
            R = min(j-1, opt[i+1][j])
            for k in range(L, R+1):
                val = dp[i][k] + dp[k+1][j] + cost
                if val<=mini:
                    mini = val
                    opt[i][j] = k
            dp[i][j] = mini
    return dp[0][n-1]

def DnC(prev, dp, cost):
    # conditions
    # opt[i][j] <= opt[i][j+1]

    # for i in range(1, K + 1):
    #	 DnC(dp_prev, dp_cur, cost, n)
    #	 dp_prev, dp_cur = dp_cur, dp_prev
    def solve(l, r, opt_l, opt_r):
        if l > r:
            return

        mid = (l + r) // 2
        best_k = -1
        best_val = float('inf')
        end = max(0,min(mid - 1, opt_r))
        for k in range(opt_l, end + 1):
            val = prev[k] + cost(k, mid)
            if val < best_val:
                best_val = val
                best_k = k
        dp[mid] = best_val
        solve(l, mid - 1, opt_l, best_k)
        solve(mid + 1, r, best_k, opt_r)
    solve(0,len(dp)-1,0,len(dp)-1)
\end{lstlisting}

Maximal Bipartite Matching
\begin{lstlisting}[language=Python]

def kuhn(d,n,m):
    # n -> left side
    # m -> right side
    # d[left side] -> [right side]
    # Maximal Bipartite matching
    # O(E V)
    match_r = [-1]*m
    match_l = [-1]*n
    for start in range(n):
        if match_l[start]!=-1:
            continue
        stack = [start]
        parent_l = [-1]*n
        parent_r = [-1]*m
        vis_r = [False]*m
        found_v = -1
        while stack:
            u = stack.pop()
            for v in d[u]:
                if vis_r[v]:
                    continue
                vis_r[v] = True
                parent_r[v] = u

                if match_r[v]==-1:
                    found_v = v
                    stack = []
                    break
                else:
                    next_u = match_r[v]
                    if parent_l[next_u]==-1:
                        parent_l[next_u] = v
                        stack.append(next_u)
        if found_v==-1:
            continue
        v = found_v
        while v!=-1:
            u = parent_r[v]
            next_v = match_l[u]
            match_l[u] = v
            match_r[v] = u
            v = next_v
    result = []
    for u in range(n):
        if match_l[u]!=-1:
            result.append((u, match_l[u]))
    return result

class HopcroftKarp:
    # n_left on left side and n_right on right side
    # O(E root(V))
    def __init__(self, n_left, n_right):
        self.n_left = n_left
        self.n_right = n_right
        self.g = [[] for _ in range(n_left)]
        self.matchL = [-1] * n_left
        self.matchR = [-1] * n_right
        self.dist = [-1] * n_left
    def add(self, u, v):
        self.g[u].append(v)
    def bfs(self):
        q = deque()
        for u in range(self.n_left):
            if self.matchL[u] == -1:
                self.dist[u] = 0
                q.append(u)
            else:
                self.dist[u] = -1
        while q:
            u = q.popleft()
            for v in self.g[u]:
                if self.matchR[v] != -1 and self.dist[self.matchR[v]] == -1:
                    self.dist[self.matchR[v]] = self.dist[u] + 1
                    q.append(self.matchR[v])
    def dfs(self, u):
        for v in self.g[u]:
            if self.matchR[v] == -1 or (self.dist[self.matchR[v]] == self.dist[u] + 1 and self.dfs(self.matchR[v])):
                self.matchL[u] = v
                self.matchR[v] = u
                return True
        return False
    def max_matching(self,flag=True):
        flow = 0
        if flag:
            for adj in self.g:
                random.shuffle(adj)
        while True:
            self.bfs()
            aug = 0
            for u in range(self.n_left):
                if self.matchL[u] == -1:
                    if self.dfs(u):
                        aug += 1
            if aug == 0:
                break
            flow += aug
        return flow
    def get_edges(self):
        return [(u, self.matchL[u]) for u in range(self.n_left) if self.matchL[u] != -1]
    def vertex_cover(self):
        visL = [False]*self.n_left
        visR = [False]*self.n_right
        q = deque()
        for u in range(self.n_left):
            if self.matchL[u]==-1:
                q.append(u)
                visL[u] = True
        while q:
            u = q.popleft()
            for v in self.g[u]:
                if not visR[v] and self.matchL[u]!=v:
                    visR[v] = True
                    if self.matchR[v]!=-1 and not visL[self.matchR[v]]:
                        visL[self.matchR[v]] = True
                        q.append(self.matchR[v])
        coverL = [u for u in range(self.n_left) if not visL[u]]
        coverR = [v for v in range(self.n_right) if visR[v]]
        return coverL,coverR
\end{lstlisting}

Graph Max Flow
\begin{lstlisting}[language=Python]
class HLPP:
    # O(V²E), still sometime unexpectedly faster than Dinic
    def __init__(self,n):
        self.n=n
        self.g=[[] for _ in range(n)]
    def add_edge(self,u,v,c):
        self.g[u].append([v,c,len(self.g[v])])
        self.g[v].append([u,0,len(self.g[u])-1])
    def maxflow(self,s,t):
        n,g=self.n,self.g
        h=[0]*n; ex=[0]*n
        seen=[0]*n
        h[s]=n
        for v,c,rev in g[s]:
            if c:
                g[v][rev][1]+=c
                ex[v]+=c
        q=deque([i for i in range(n) if i!=s and i!=t and ex[i]>0])
        def push(u,e):
            v,c,rev=e
            d=min(ex[u],c)
            if h[u]==h[v]+1 and d>0:
                e[1]-=d; g[v][rev][1]+=d
                ex[u]-=d; ex[v]+=d
                if ex[v]==d and v!=s and v!=t:
                    q.append(v)
                return True
            return False
        def relabel(u):
            h[u]=min(h[v]+1 for v,c,_ in g[u] if c>0)
        while q:
            u=q[0]
            while ex[u]:
                if seen[u]<len(g[u]):
                    if not push(u,g[u][seen[u]]):
                        seen[u]+=1
                else:
                    relabel(u)
                    seen[u]=0
            q.popleft()
        return ex[t]

class Dinic:
    # O(V²E)
    # Unit capacities/bipartite/DAG: ~O(E√V)
    def __init__(self, n):
        self.lvl = [0] * n
        self.ptr = [0] * n
        self.q = [0] * n
        self.adj = [[] for _ in range(n)]
    def add_edge(self, a, b, c, rcap=0):
        self.adj[a].append([b, len(self.adj[b]), c, 0])
        self.adj[b].append([a, len(self.adj[a]) - 1, rcap, 0])
    def dfs(self, v, t, f):
        if v == t or not f:
            return f
        for i in range(self.ptr[v], len(self.adj[v])):
            e = self.adj[v][i]
            if self.lvl[e[0]] == self.lvl[v] + 1:
                p = self.dfs(e[0], t, min(f, e[2] - e[3]))
                if p:
                    self.adj[v][i][3] += p
                    self.adj[e[0]][e[1]][3] -= p
                    return p
            self.ptr[v] += 1
        return 0
    def calc(self, s, t):
        INF = float('inf')
        flow, self.q[0] = 0, s
        for l in range(31):  # l = 30 maybe faster for random data
            while True:
                self.lvl, self.ptr = [0] * len(self.q), [0] * len(self.q)
                qi, qe, self.lvl[s] = 0, 1, 1
                while qi < qe and not self.lvl[t]:
                    v = self.q[qi]
                    qi += 1
                    for e in self.adj[v]:
                        if not self.lvl[e[0]] and (e[2] - e[3]) >> (30 - l):
                            self.q[qe] = e[0]
                            qe += 1
                            self.lvl[e[0]] = self.lvl[v] + 1
                p = self.dfs(s, t, INF)
                while p:
                    flow += p
                    p = self.dfs(s, t, INF)
                if not self.lvl[t]:
                    break
        return flow
    def get_paths(self,s,t):
        # returns that paths taken for the maxflow
        paths=[]
        while True:
            stack=[(s,[],float('inf'),set([s]))]
            found=False
            while stack:
                v,path,f,vis=stack.pop()
                if v==t:
                    pushed=f
                    res_path=path
                    found=True
                    break
                for i,e in enumerate(self.adj[v]):
                    to,rev,cap,flow=e
                    if flow>0 and to not in vis:
                        stack.append((to,path+[(v,to)],min(f,flow),vis|{to}))
            if not found:
                break
            f=pushed
            for u,w in res_path:
                for e in self.adj[u]:
                    if e[0]==w and e[3]>0:
                        e[3]-=f
                        self.adj[w][e[1]][3]+=f
                        break
            paths.append((pushed,res_path))
        return paths
\end{lstlisting}

Link Cut Tree
\begin{lstlisting}[language=Python]

class LinkCutTree:
    class Node:
        __slots__=('l','r','p','rev','val','sm')
        def __init__(self,val=0):
            self.l=-1
            self.r=-1
            self.p=-1
            self.rev=False
            self.val=val
            self.sm=val
    def __init__(self,n):
        self.n=n
        self.t=[self.Node() for _ in range(n)]
    def _is_root(self,x):
        p=self.t[x].p
        return p==-1 or (self.t[p].l!=x and self.t[p].r!=x)
    def _push(self,x):
        if x==-1:
            return
        t=self.t[x]
        if t.rev:
            t.l,t.r=t.r,t.l
            if t.l!=-1:
                self.t[t.l].rev^=1
            if t.r!=-1:
                self.t[t.r].rev^=1
            t.rev=False
    def _pull(self,x):
        t=self.t[x]
        t.sm=t.val
        if t.l!=-1:
            t.sm+=self.t[t.l].sm
        if t.r!=-1:
            t.sm+=self.t[t.r].sm
    def _rotate(self,x):
        p=self.t[x].p
        g=self.t[p].p
        if self.t[p].l==x:
            b=self.t[x].r
            self.t[x].r=p
            self.t[p].l=b
            if b!=-1:
                self.t[b].p=p
        else:
            b=self.t[x].l
            self.t[x].l=p
            self.t[p].r=b
            if b!=-1:
                self.t[b].p=p
        self.t[p].p=x
        self.t[x].p=g
        if g!=-1:
            if self.t[g].l==p:
                self.t[g].l=x
            elif self.t[g].r==p:
                self.t[g].r=x
        self._pull(p)
        self._pull(x)
    def _splay(self,x):
        stk=[]
        y=x
        while True:
            stk.append(y)
            if self._is_root(y):
                break
            y=self.t[y].p
        while stk:
            self._push(stk.pop())
        while not self._is_root(x):
            p=self.t[x].p
            g=self.t[p].p
            if not self._is_root(p):
                if (self.t[g].l==p)==(self.t[p].l==x):
                    self._rotate(p)
                else:
                    self._rotate(x)
            self._rotate(x)
    def _access(self,x):
        last=-1
        while x!=-1:
            self._splay(x)
            self.t[x].r=last
            self._pull(x)
            last=x
            x=self.t[x].p
    def makeroot(self,x):
        self._access(x)
        self._splay(x)
        self.t[x].rev^=1
    def findroot(self,x):
        self._access(x)
        self._splay(x)
        while True:
            self._push(x)
            if self.t[x].l==-1:
                break
            x=self.t[x].l
        self._splay(x)
        return x
    def connected(self,u,v):
        return self.findroot(u)==self.findroot(v)
    def link(self,u,v):
        self.makeroot(u)
        if self.findroot(v)==u:
            return False
        self.t[u].p=v
        return True
    def cut(self,u,v):
        self.makeroot(u)
        self._access(v)
        self._splay(v)
        if self.t[v].l!=u or self.t[u].r!=-1:
            return False
        self.t[u].p=-1
        self.t[v].l=-1
        self._pull(v)
        return True
    def set_value(self,u,val):
        self._access(u)
        self._splay(u)
        self.t[u].val=val
        self._pull(u)
    def path_sum(self,u,v):
        self.makeroot(u)
        self._access(v)
        self._splay(v)
        return self.t[v].sm
\end{lstlisting}

Centroid Decomposition
\begin{lstlisting}[language=Python]

def shallowest_decomposition_tree(graph, root=0):
    # returns a root and a directed tree, iterate on it and you
    # will find the new centroid
    # slightly faster than centroid_decomposition_tree
    ctz = lambda x:(x & -x).bit_length()-1
    n = len(graph)
    forbidden = [0] * n
    decomposition_tree = [[] for _ in range(n)]
    stacks = [[] for _ in range(n.bit_length())]
    def extract_chain(labels, u):
        while labels:
            label = labels.bit_length() - 1
            labels ^= 2**label
            v = stacks[label].pop()
            decomposition_tree[u].append(v)
            u = v
    dfs = [root]
    while dfs:
        u = dfs.pop()
        if u >= 0:
            forbidden[u] = -1
            dfs.append(~u)
            for v in graph[u]:
                if not forbidden[v]:
                    dfs.append(v)
        else:
            u = ~u
            forbidden_once = forbidden_twice = 0
            for v in graph[u]:
                forbidden_twice  |= forbidden_once & (forbidden[v] + 1)
                forbidden_once  |= forbidden[v] + 1
            forbidden[u] = forbidden_once | (2**forbidden_twice.bit_length() - 1)
            label_u = ctz(forbidden[u] + 1)
            stacks[label_u].append(u)
            for v in graph[u]:
                extract_chain((forbidden[v] + 1) & (2**label_u - 1), u)
    max_label = (forbidden[root] + 1).bit_length() - 1
    decomposition_root = stacks[max_label].pop()
    extract_chain((forbidden[root] + 1) & (2**max_label - 1), decomposition_root)
    return decomposition_tree, decomposition_root

def centroid_decomposition_tree(graph):
    # returns a root and a directed tree, iterate on it and you
    # will find the new centroid
    n = len(graph)
    graph = [c[:] for c in graph]
    bfs = [0]
    for node in bfs:
        bfs += graph[node]
        for nei in graph[node]:
            graph[nei].remove(node)
    size = [0] * n
    for node in reversed(bfs):
        size[node] = 1 + sum(size[child] for child in graph[node])
    decomposition_tree = [[] for _ in range(n)]
    def centroid_reroot(u):
        N = size[u]
        while True:
            for v in graph[u]:
                if size[v] > N // 2:
                    size[u] = N - size[v]
                    graph[u].remove(v)
                    graph[v].append(u)
                    u = v
                    break
            else:
                decomposition_tree[u] = [centroid_reroot(v) for v in graph[u]]
                return u
    decomposition_root = centroid_reroot(0)
    return decomposition_tree, decomposition_root
\end{lstlisting}

Lazy Segment Tree
\begin{lstlisting}[language=Python]

class SegmentTree:
    def __init__(self, data, default=0, mode='s',func=lambda x,y:x+y):
        self.func = func
        self.mode = mode
        self._default = default
        self.n = len(data)
        self.size = 1 << (self.n - 1).bit_length()
        self.tree = [default] * (2 * self.size)
        self._size = [0] * (2 * self.size)
        self._size[self.size:] = [1] * self.size
        for i in range(self.size - 1, 0, -1):
            self._size[i] = self._size[i << 1] + self._size[i << 1 | 1]
        self.lazy_add = 0
        self.lazy_set = None
        self.lazy_add = [0] * self.size
        self.lazy_set = [None] * self.size
        for i in range(self.n):
            self.tree[self.size + i] = data[i]
        for i in range(self.size - 1, 0, -1):
            self.tree[i] = self.func(self.tree[i << 1], self.tree[i << 1 | 1])
    def _apply_set(self, pos, value):
        if self.mode == 's':
            self.tree[pos] = value * self._size[pos]
        else:
            self.tree[pos] = value
        if pos < self.size:
            self.lazy_set[pos] = value
            self.lazy_add[pos] = 0
    def _apply_add(self, pos, value):
        if self.mode == 's':
            self.tree[pos] += value * self._size[pos]
        else:
            self.tree[pos] += value
        if pos < self.size:
            if self.lazy_set[pos] is not None:
                self.lazy_set[pos] += value
            else:
                self.lazy_add[pos] += value
    def _build(self, pos):
        while pos > 1:
            pos >>= 1
            self.tree[pos] = self.func(self.tree[pos << 1], self.tree[pos << 1 | 1])
            if self.lazy_set[pos] is not None:
                if self.mode == 's':
                    self.tree[pos] = self.lazy_set[pos] * self._size[pos]
                else:
                    self.tree[pos] = self.lazy_set[pos]
            if self.lazy_add[pos] != 0:
                if self.mode == 's':
                    self.tree[pos] += self.lazy_add[pos] * self._size[pos]
                else:
                    self.tree[pos] += self.lazy_add[pos]
    def _push(self, pos):
        for shift in range(self.size.bit_length() - 1, 0, -1):
            i = pos >> shift
            set_val = self.lazy_set[i]
            if set_val is not None:
                self._apply_set(i << 1, set_val)
                self._apply_set(i << 1 | 1, set_val)
                self.lazy_set[i] = None
            add_val = self.lazy_add[i]
            if add_val != 0:
                self._apply_add(i << 1, add_val)
                self._apply_add(i << 1 | 1, add_val)
                self.lazy_add[i] = 0
    def range_update(self, left, right, value,flag=True):
        # Range Update in [L,R] if flag, then add
        if flag:
            l = left + self.size
            r = right + self.size
            l0, r0 = l, r
            self._push(l0)
            self._push(r0)
            while l <= r:
                if l & 1: self._apply_add(l, value); l += 1
                if not r & 1: self._apply_add(r, value); r -= 1
                l >>= 1; r >>= 1
            self._build(l0)
            self._build(r0)
        else:
            l = left + self.size
            r = right + self.size
            l0, r0 = l, r
            self._push(l0)
            self._push(r0)
            while l <= r:
                if l & 1: self._apply_set(l, value); l += 1
                if not r & 1: self._apply_set(r, value); r -= 1
                l >>= 1; r >>= 1
            self._build(l0)
            self._build(r0)
    def range_query(self, left, right):
        # Range Query in [L,R]
        l = left + self.size
        r = right + self.size
        self._push(l)
        self._push(r)
        res = self._default
        while l <= r:
            if l & 1: res = self.func(res, self.tree[l]); l += 1
            if not r & 1: res = self.func(res, self.tree[r]); r -= 1
            l >>= 1; r >>= 1
        return res
    def __repr__(self):
        return f"SegmentTree({[self.range_query(i,i) for i in range(self.n)]})"
\end{lstlisting}

Utils
\begin{lstlisting}[language=Python]
getcontext().prec = 50;sys.setrecursionlimit(10**6)
sys.set_int_max_str_digits(10**5)
import builtins
globals()['print'] = lambda *args, **kwargs: builtins.print(*args, flush=True, **kwargs)
\end{lstlisting}

\section{Definitions}
\begin{description}
\item[Extended GCD] Computes $x,y$ s.t. $ax+by=\gcd(a,b)$.
\item[CRT] Finds smallest $a$ satisfying $a\equiv b_i \pmod{m_i}$.
\item[$\phi(n)$] Count of integers $<n$ coprime with $n$.
\item[Suffix Array] Starting indices of suffixes sorted lexicographically.
\item[Euler Path] Uses every edge exactly once.
\item[Bridge] Edge whose removal increases components.
\item[Articulation Point] Vertex whose removal increases components.
\item[2ECC] Maximal subgraph where every pair lies on a cycle.
\end{description}

\section{C++ Notes}

Use \texttt{get<i>(tuple)} for tuple access.

Custom sort:
\texttt{[](auto\& a,auto\& b)\{return get<2>(a)<get<2>(b);\}}

2D vector init:
\texttt{vector<vector<int>> dp(n, vector<int>(m,0));}

Prefer \texttt{emplace\_back} for complex objects.

Min-heap:
\texttt{priority\_queue<pi,vector<pi>,greater<pi>> pq;}


\section{CP - C++}
\begin{lstlisting}[language=C++]
#include <bits/stdc++.h>
#include <ext/pb_ds/assoc_container.hpp>
#include <ext/pb_ds/tree_policy.hpp>
using namespace __gnu_pbds;
using namespace std;

#define int long long
#define INF LLONG_MAX
#define fastio ios::sync_with_stdio(false); cin.tie(0);
#define len(arr) arr.size()
#define f first
#define s second
#define pb push_back
using vi = vector<int>;
using vii = vector<vector<int>>;
using pi = pair<int,int>;
#pragma GCC optimize("O3,unroll-loops")
#pragma GCC target("avx2,bmi,bmi2,lzcnt,popcnt")

template <typename T>
using ordered_multiset = tree<
    std::pair<T, int>,
    null_type,
    std::less<std::pair<T, int>>,
    rb_tree_tag,
    tree_order_statistics_node_update
>;

int II(){int a;cin>>a;return a;}
string SI(){string s;cin>>s;return s;}
vi LII(int n){vi a(n);cin>>a;return a;}

void solve() {

}

int32_t main() {
    fastio
    int t=1;
    cin >> t;
    while (t--) solve();
    return 0;
}
\end{lstlisting}


\section{Sorted List - C++}
\begin{lstlisting}[language=C++]
template <typename T>
class SortedList {
    ordered_multiset<T> os;
    int uid = 0;
public:
    void insert(T x) {
        os.insert({x, uid++});
    }
    void erase_one(T x) {
        auto it = os.lower_bound({x, 0});
        if (it != os.end() && it->first == x)
            os.erase(it);
    }
    int index(T x) {
        return os.order_of_key({x, 0});
    }
    int count(T x) {
        return os.order_of_key({x, INT_MAX}) - os.order_of_key({x, 0});
    }
    T operator[](int k) {
        return os.find_by_order(k)->first;
    }
    int size() const {
        return os.size();
    }
    void print() {
        for (int i = 0; i < size(); i++) std::cout << (*this)[i] << " ";
        std::cout << "\n";
    }
};
\end{lstlisting}

\section{Lazy Segment Tree - C++}
\begin{lstlisting}[language=C++]
class SegmentTree {
public:
    int n;
    vector<int> tree;
    vector<int> lazy_add;
    vector<int> lazy_set;
    int NO_ASSIGNMENT = LLONG_MIN;
    SegmentTree(vector<int>& arr) {
        n = arr.size();
        tree.assign(4 * n, 0);
        lazy_add.assign(4 * n, 0);
        lazy_set.assign(4 * n, NO_ASSIGNMENT);
        build_tree(1, 0, n - 1, arr);
    }
    static int func(int a, int b) {
        return a + b;
    }
    void build_tree(int node, int start, int end, vector<int>& arr) {
        if (start == end) {
            tree[node] = arr[start];
        } else {
            int mid = start + (end - start) / 2;
            build_tree(2 * node, start, mid, arr);
            build_tree(2 * node + 1, mid + 1, end, arr);
            tree[node] = func(tree[2 * node], tree[2 * node + 1]);
        }
    }
    void propagate_lazy(int node, int start, int end) {
        int current_range_size = (end - start + 1);
        if (lazy_set[node] != NO_ASSIGNMENT) {
            tree[node] = lazy_set[node] * current_range_size;
            if (start != end) {
                lazy_set[2 * node] = lazy_set[node];
                lazy_set[2 * node + 1] = lazy_set[node];
                lazy_add[2 * node] = 0;
                lazy_add[2 * node + 1] = 0;
            }
            lazy_set[node] = NO_ASSIGNMENT;
        }
        if (lazy_add[node] != 0) {
            tree[node] += lazy_add[node] * current_range_size;
            if (start != end) {
                lazy_add[2 * node] += lazy_add[node];
                lazy_add[2 * node + 1] += lazy_add[node];
            }
            lazy_add[node] = 0;
        }
    }
    void update(int node, int start, int end, int l, int r, int value, bool is_add) {
        propagate_lazy(node, start, end);
        if (start > r || end < l) {
            return;
        }
        if (start >= l && end <= r) {
            if (is_add) {
                tree[node] += value * (end - start + 1);
                if (start != end) {
                    lazy_add[2 * node] += value;
                    lazy_add[2 * node + 1] += value;
                }
            } else {
                tree[node] = value * (end - start + 1);
                if (start != end) {
                    lazy_set[2 * node] = value;
                    lazy_set[2 * node + 1] = value;
                    lazy_add[2 * node] = 0;
                    lazy_add[2 * node + 1] = 0;
                }
            }
            return;
        }
        int mid = start + (end - start) / 2;
        update(2 * node, start, mid, l, r, value, is_add);
        update(2 * node + 1, mid + 1, end, l, r, value, is_add);
        tree[node] = func(tree[2 * node], tree[2 * node + 1]);
    }
    int query(int node, int start, int end, int l, int r) {
        propagate_lazy(node, start, end);
        if (start > r || end < l) {
            return 0;
        }
        if (start >= l && end <= r) {
            return tree[node];
        }
        int mid = start + (end - start) / 2;
        return func(query(2 * node, start, mid, l, r), query(2 * node + 1, mid + 1, end, l, r));
    }
    void range_update(int l, int r, int value, bool is_add = true) {
        update(1, 0, n - 1, l, r, value, is_add);
    }
    int range_query(int l, int r) {
        return query(1, 0, n - 1, l, r);
    }
    vector<int> to_list() {
        vector<int> result(n);
        for (int i = 0; i < n; ++i) {
            result[i] = range_query(i, i);
        }
        return result;
    }
};
\end{lstlisting}

\section{FFT - C++}
\begin{lstlisting}[language=C++]
class FFT {
public:
    static const int MOD = 998244353;
    static const int MOD1 = 469762049;
    int MOD2, mod_inv;
    vector<int> W, W1, iW, iW1;
    FFT() {
        MOD2 = power(MOD, MOD1 - 2, MOD1);
        mod_inv = XT_GCD(MOD, MOD1).second % MOD1;
        int g = 3;
        int ig = power(g, MOD - 2, MOD);
        int ig1 = power(g, MOD1 - 2, MOD1);
        W.resize(30); W1.resize(30); iW.resize(30); iW1.resize(30);
        for (int i = 0; i < 30; ++i) {
            W[i] = power(g, (MOD - 1) >> i, MOD);
            W1[i] = power(g, (MOD1 - 1) >> i, MOD1);
            iW[i] = power(ig, (MOD - 1) >> i, MOD);
            iW1[i] = power(ig1, (MOD1 - 1) >> i, MOD1);
        }
    }
    int power(int base, int exp, int mod) {
        int res = 1;
        while (exp) {
            if (exp % 2) res = res * base % mod;
            base = base * base % mod;
            exp /= 2;
        }
        return res;
    }
    pair<int, int> XT_GCD(int a, int b) {
        if (b == 0) return {a, 1};
        pair<int,int> k9 = XT_GCD(b, a % b);
        int g = k9.first;int x1 = k9.second;
        return {g, x1 - (a / b) * x1};
    }

    int CRT(int a, int mod1, int b, int mod2) {
        return (a + (b - a) * mod_inv % mod2 * mod1) % (mod1 * mod2);
    }
    void fft(int k, vector<int>& f, vector<int>& f1) {
        for (int l = k; l > 0; --l) {
            int d = 1 << (l - 1);
            vector<pair<int, int>> U = {{1, 1}};
            for (int i = 0; i < d; ++i)
                U.emplace_back(U.back().first * W[l] % MOD, U.back().second * W1[l] % MOD1);
            for (int i = 0; i < (1 << (k - l)); ++i) {
                for (int j = 0; j < d; ++j) {
                    int s = i * 2 * d + j;
                    int tmp_f = f[s] - f[s + d];
                    int tmp_f1 = f1[s] - f1[s + d];
                    f[s] = (f[s] + f[s + d]) % MOD;
                    f[s + d] = U[j].first * tmp_f % MOD;
                    f1[s] = (f1[s] + f1[s + d]) % MOD1;
                    f1[s + d] = U[j].second * tmp_f1 % MOD1;
                }
            }
        }
    }
    void ifft(int k, vector<int>& f, vector<int>& f1) {
        for (int l = 1; l <= k; ++l) {
            int d = 1 << (l - 1);
            for (int i = 0; i < (1 << (k - l)); ++i) {
                int u = 1, u1 = 1;
                for (int j = i * 2 * d; j < (i * 2 + 1) * d; ++j) {
                    f[j + d] = f[j + d] * u % MOD;
                    f[j] = (f[j] + f[j + d]) % MOD;
                    f[j + d] = (f[j] - f[j + d] + MOD) % MOD;
                    u = u * iW[l] % MOD;
                    f1[j + d] = f1[j + d] * u1 % MOD1;
                    f1[j] = (f1[j] + f1[j + d]) % MOD1;
                    f1[j + d] = (f1[j] - f1[j + d] + MOD1) % MOD1;
                    u1 = u1 * iW1[l] % MOD1;
                }
            }
        }
    }
    vector<int> convolve(vector<int> A, vector<int> B) {
        int n0 = A.size() + B.size() - 1;
        int k = 0;
        while ((1 << k) < n0) ++k;
        int n = 1 << k;
        A.resize(n, 0);
        B.resize(n, 0);
        vector<int> A1 = A, B1 = B;
        fft(k, A, A1);
        fft(k, B, B1);
        for (int i = 0; i < n; ++i) {
            A[i] = A[i] * B[i] % MOD;
            A1[i] = A1[i] * B1[i] % MOD1;
        }
        ifft(k, A, A1);
        int inv = power(n, MOD - 2, MOD);
        int inv1 = power(n, MOD1 - 2, MOD1);
        A.resize(n0);
        for (int i = 0; i < n0; ++i) {
            A[i] = CRT(A[i] * inv % MOD, MOD, A1[i] * inv1 % MOD1, MOD1);
        }
        return A;
    }
};
\end{lstlisting}

\section{I/O - Shubham}
\begin{lstlisting}[language=C++]
#define all(a) (a).begin(),(a).end()
#define pb push_back
#define sz(a) ((int)(a).size())
#define mp make_pair
#define fi first
#define se second

typedef pair<int,int> pint;
typedef long long ll;
typedef unsigned long long ull;
typedef vector<int> vi;
sort(v.begin(), v.end());
v.resize(unique(v.begin(), v.end()) - v.begin());
__builtin_popcount(unsigned int)->no of set bits
__builtin_ffs(int)->most right set bit(lsb)
memset(dp, -1, sizeof(dp));

     int n;
     cin >> n;
     int arr[n];
     for(int i=0;i<n;i++){
	 cin >> arr[i];
      }

 bitset<8> decimalBitset(15);
 to_ullong()
 to_string()
 flip()
 set()
 reset()
 count()
 bitset<4> allSet("1111"), allUnset;

 *min_element(a.begin(), a.end());

#define int long long
#define fr(i, a, b) for (int i = (a); i < (int)(b); i++)
is_sorted(A, A + range1)
count(arr, arr + n, 3);
int n = sizeof(arr) / sizeof(arr[0]);



ios_base::sync_with_stdio(false);
cin.tie(NULL);

map<int, int>::iterator it = myMap.find(key);
int position = distance(myMap.begin(), it);


      int arr[n];
      map<int,int,greater<int>> m;
       for(int i=0;i<=30;i++) m[i]=0;
       for(int i=0;i<n;i++){
        cin >> arr[i];
        for(int j=0;j<31;j++){
            if((arr[i]>>j)&1!=0) m[j]++;
            }
 
       }


 int n = stoi(s);

 template <class T1, class T2, class Pred = std::less<T2> >
struct sort_pair_second {
    bool operator()(const std::pair<T1,T2>&left, const std::pair<T1,T2>&right) {
        Pred p;
        return p(left.second, right.second);
    }
};


std::sort(v.begin(), v.end(), sort_pair_second<int, int>());

std::sort(v.begin(), v.end(), sort_pair_second<int, int, std::greater<int> >());



int n,e;
cin >> n >> e;
while(e--){
    int x,y;
    cin >> x >> y;
    g[x].push_back(y);
    g[y].push_back(x);
}
\end{lstlisting}

\section{DFS - Shubham}
\begin{lstlisting}[language=C++]
// for dfs
const int N=1e5+10;
vector<int> g[N];
bool check[N];
void dfs(int vertex){
    check[vertex]=true;
    for(auto child: g[vertex]){
        if(check[child]) continue;
        
        dfs(child);
    }
    return;
}
\end{lstlisting}

\section{Maths - Shubham}
\begin{lstlisting}[language=C++]
vector<int> prime;
for(int i=2;i*i<=n;i++){
    while(n%i==0){
    prime.push_back(i);
    n/=i;
    }
}
if(n>1) prime.push_back(n);
int ans=LLONG_MAX;       

vector<int> primenums(int n) {
vector<int> isprime(n + 1, 1);
isprime[0] = isprime[1] = 0;

for (int i = 2; i * i <= n; ++i) {
    if (isprime[i]) {
        for (int j = i * i; j <= n; j += i) {
            isprime[j] = 0;
        }
    }
}

vector<int> primes;
for (int i = 2; i <= n; ++i) {
    if (isprime[i]) {
        primes.push_back(i);
    }
}

return primes;
}

int n = 100; 
vector<int> primes=primenums(n);

vector<int> getFactors(int n) {
    vector<int> factors;
    for(int i = 1; i <= sqrt(n); i++) {
        if(n % i == 0) {
            factors.push_back(i);
            if(i != n / i) {
                factors.push_back(n / i);
            }
        }
    }
    sort(factors.begin(), factors.end());
    return factors;
}
\end{lstlisting}

\section{Strings - Shubham}
\begin{lstlisting}[language=C++]
vector<int> rabin_karp(string const& s, string const& t) {
    const int p = 31; 
    const int m = 1e9 + 9;
    int S = s.size(), T = t.size();

    vector<long long> p_pow(max(S, T)); 
    p_pow[0] = 1; 
    for (int i = 1; i < (int)p_pow.size(); i++) 
        p_pow[i] = (p_pow[i-1] * p) % m;

    vector<long long> h(T + 1, 0); 
    for (int i = 0; i < T; i++)
        h[i+1] = (h[i] + (t[i] - 'a' + 1) * p_pow[i]) % m; 
    long long h_s = 0; 
    for (int i = 0; i < S; i++) 
        h_s = (h_s + (s[i] - 'a' + 1) * p_pow[i]) % m; 

    vector<int> occurrences;
    for (int i = 0; i + S - 1 < T; i++) {
        long long cur_h = (h[i+S] + m - h[i]) % m;
        if (cur_h == h_s * p_pow[i] % m)
            occurrences.push_back(i);
    }
    return occurrences;
}

class Solution {
public:
    int M=1e9+7;
    int strStr(string haystack, string needle) {
        int n=haystack.size();
        int m=needle.size();
        if(m>n) return -1;
        map<long long,int> mp;
        int pm=31;
        vector<long long> p_pow(n+1);
        p_pow[0]=1;
        for(int i=1;i<=n;i++){
            p_pow[i]=(p_pow[i-1]*pm)%M;            
        }
        vector<long long> h(n+1,0);
        h[1]=(haystack[0]-'a'+1);
        for(int i=0;i<n;i++){
            h[i+1]=(h[i]+(haystack[i]-'a'+1)*p_pow[i])%M;
        }

        int hash=0;
        for(int i=0;i<m;i++){
            hash+=((needle[i]-'a'+1)*(p_pow[i]))%M;
            hash%=M;
        }

        for(int i=0;i<(n-m+1);i++){
            int curr=(h[i+m]-h[i]+M)%M;
            if(curr==(hash*p_pow[i])%M) return i;
        }
        return -1;
    }
}; 
\end{lstlisting}

\section{Modular Exponentation - Shubham}
\begin{lstlisting}[language=C++]
int modularExponentiation(int base, int exp, int mod) {
    int result = 1;
    base = base % mod;
    while (exp > 0) {
        if (exp % 2 == 1) {
            result = (result * base) % mod;
        }
        exp = exp >> 1;
        base = (base * base) % mod;
    }
    return result;
}
\end{lstlisting}

\section{Usage}
\begin{lstlisting}[]
- **Number Theory 1(numtheory)**: 
    - **miller_isprime**: Calculates if a number is prime or not in 7log2n, better for bigger numbers.
- **Number Theory 2(numrare)**:
    - **pollard_rho**: Returns a random factor of n.
    - **miller_pollard_rho**: Miller along with combination of pollard rho.
    - **prime_factors_large**: Probablistic algo, returns prime factors in n^(1/4).
    - **Discrete Log**: returns smallest x such that pow(a,x,mod)=b
    - **composite_crt**: Chinese Remainder Theorem. **
    - **phi**: Computes Phi(n) for all n. **
- **Linear Algebra**:
    - **compute_rank**: Computes rank of a matrix using Gaussian Elimination.
    - **Xor Basis**: Returns array of elements such that any element added to the data structure can be made using their combination.
    - **Simplex Bland**: Solves LPP Problems using Simplex Method.
    - **s_lr**: Computer Summation of i**power for l<=i<=r, given power<=3
- **Combinatorics(pnc)**:
    - **catalan**: (1/(n+1)) * (2nCn)
    - **derangements**: Number of permutations of n objects where no object appears in its original position
    - **stirling_2**: Number of ways to partition n elements into k non-empty subsets
    - **partition**: Ways to partition n into k or fewer parts of size 1 or greater
    - **Lucas nCk**: Used to find nCk for huge n and k modulo p in log p(n)
    - **beggar**: number of ways to choose k elements >= 0 such that their sum is n
- **Geometry**:
    - **onSegment**: returns True if q lies on pr
    - **orientation**: Checks orientation of p,q,r
    - **doIntersect**: Checks if two line segments intersect
    - **get_2dline**: Returns normalized 2d line
    - **is_parallel**: Checks if lines are parallel
    - **is_same**: Checks if two lines are same
    - **collinear**: Checks if three points are collinear
    - **intersect**: Returns intersection point of 2 lines
    - **rotate**: Rotates point p around origin by thetha(radians)
    - **perimeter**: Gives perimeter of a polygon
    - **area**: Polygon area (shoelace formula)
    - **centroid**: Returns the centroid of a polygon
    - **is_in_circle**: checks if point p is inside circle centered at c with radius r
    - **incircle_radius**: Radius incircle of triangle
    - **circumecircle_radius**: Radius of circumcircle
    - **Convex Hull**: Returns minimal set of points, within which all other points are contained
- **String Functions(sf)**:
    - **KMP**: Longest Prefix Matching ending at position i.
    - **Z-function**: Longest Prefix Matching starting at position i.
    - **Manacher's Algorithm**: Longest Palindrome.
    - **Suffix Array**: Suffix Array. **
        - **SuffArr**: O(N) theoretical complexity but slower
        - **SuffArray**: O(NlogN) but faster, due to less constant time
    - **LCP**: Longest common prefix between adjacent elements in Suffix Array.
- **Graph 8(graph_decom)**: Centroid Decompositon (Similar to Merge Sort)
    - **Shallowest Decomposition Tree**: O(N) Decomposition
    - **Centroid Decomposition Tree**: O(N log N) Standard Decomposition
- **Graph 9(graphflow)**: Graph Flow
    - **Dinic**: Maximal graph flow V2E
    - **HLPP**: Maximal Graph flow V2root(E)
- **Graph 10(graphmatch)**: Maximal Bipartite Matching
    - **Kuhn's Algorithm**: E*V
    - **Hopcroft Karp Algorithm**: E*root(V) - Maximal Bipartite Matching -> return min vertex cover/ max edges
- **Graph 11(lct)**: LCT
    - **Link Cut Tree**: O(log n) complexity to remove, create edge, check connectivity, distance, need to add more such functions
- **DP Optimizations**:
    - **Knuth's Optimization**: 
    
        dp[i][j] = min(dp[i][k]+dp[k+1][j] over all k)+cost(i,j)
    
        Condition: 
            
            For all a<=b<=c<=d:
                cost(a,c)+cost(b,d)<=cost(a,d)+cost(b,c)
            opt[l][r-1] <= opt[l][r] <= opt[l+1][r]
    
    - **Divide and Conquer Optimization**:
        
        dp[i][j] = min(dp[l][j-1]+cost(l,i) over all l)
        
        Condition:
        
            opt[i][j] <= opt[i][j+1]
- **FWHT**: Convolution such that ans[i] = summation of a[j]*b[k], where j logical condition k = i
- **Min Cost Max Flow**: Minimum cost that can be used while maximizing the flow
    - **MCMF**: To be used when the flow in between edges is non constant, O(F E V)
    - **Hungarian**: To be used when flow is constant, and you have a cost matrix O(n3)
\end{lstlisting}

\scriptsize
\end{multicols*}
\end{document}
